\documentclass[11pt,a4paper]{article}
\usepackage[T1]{fontenc}
\usepackage[utf8]{inputenc}
\usepackage{lmodern}
\usepackage[a4paper,margin=24mm,headheight=15pt]{geometry}
\usepackage{amsmath,amssymb,amsthm}
\usepackage{booktabs,longtable,tabularx,array}
\usepackage{xcolor,graphicx,microtype}
\usepackage{xurl}
\usepackage{enumitem}
\usepackage{listings}
\usepackage{fancyhdr}
\usepackage[colorlinks=true,linkcolor=blue!45!black,urlcolor=blue!50!black,citecolor=blue!45!black]{hyperref}
\hypersetup{pdftitle={New 194-state machines for the PA consistency program},pdfauthor={Research prepared for Vladimir Reshetnikov},pdfsubject={Exact tables, iPRA consistency equivalence, census and bounded reduction searches}}
\definecolor{ink}{RGB}{22,42,62}
\definecolor{pale}{RGB}{242,246,249}
\pagestyle{fancy}
\fancyhf{}
\fancyhead[L]{\small\color{ink}PA consistency machines}
\fancyhead[R]{\small 14 November 2026}
\fancyfoot[C]{\thepage}
\setlength{\parskip}{4pt}
\setlength{\parindent}{0pt}
\setlength{\emergencystretch}{2em}
\setlist{nosep,topsep=5pt,leftmargin=*}
\lstset{basicstyle=\ttfamily\small,columns=fullflexible,keepspaces=true,breaklines=true,frame=single,rulecolor=\color{black!15},backgroundcolor=\color{pale}}
\newcommand{\ipra}{\mathsf{iPRA}}
\newcommand{\pa}{\mathsf{PA}}
\newcommand{\Con}{\operatorname{Con}}
\newcommand{\Halt}{\operatorname{Halt}}
\newcommand{\ML}{M_{\mathrm{LC}}}
\newcommand{\MLH}{M_{\mathrm{LC\_AX11H}}}
\newcommand{\snap}{32d0d8b7d93a63ad4106c18cdb29b37927acdad2}
\newcommand{\repourl}{https://github.com/VladimirReshetnikov/ProveIt/tree/32d0d8b7d93a63ad4106c18cdb29b37927acdad2}
\newtheorem{theorem}{Theorem}[section]
\newtheorem{lemma}[theorem]{Lemma}
\newtheorem{proposition}[theorem]{Proposition}

\begin{document}
\hypersetup{pageanchor=false}
\begin{titlepage}
\color{ink}
\vspace*{14mm}
{\Large Research report}\par
\vspace{12mm}
{\huge\bfseries New 194-state machines for the PA consistency program\par}
\vspace{7mm}
{\Large Comparison, iPRA proofs, and reduction searches\par}
\vspace{15mm}
{\large Prepared for Vladimir Reshetnikov\\14 November 2026}\par
\vspace{15mm}
\begin{minipage}{0.95\textwidth}
\color{black}
\textbf{Principal construction.} PA194LC\_AX11H has 194 working states and
differs from PA194AX11H in exactly three native transition rows. Together
with four new implementations of the original primary source, it supplies
five structurally new graphs relative to the 110 complete 194-state
graphs in the final audited snapshot.

\textbf{Arithmetic result.} A written proof of relative iPRA equivalence
between the new machine and AX11H gives the explicit two-way witness bound
\[
\Theta(t)=13t^2+707t+58\,616.
\]
Importing the revised AX11H packet's written absolute calibration yields
\[
\ipra\vdash\neg\Halt(\MLH)\leftrightarrow\Con(\pa).
\]

\textbf{Proof boundary.} No verified machine below 194 states was found.
The absolute corollary imports the newly supplied written calibration,
whose substantive bridge arguments were reviewed. It is not a new
end-to-end Lean or Rocq theorem. The original-primary loop family retains
its separately stated relative proof status.
\end{minipage}
\vfill
\small Initial snapshot: \path{32d0d8b7d93a63ad4106c18cdb29b37927acdad2}.\\
Final delta audit: \path{603c58ca3bd6bf3222fcfc4e308e6b079f65d573}.\\
Companion archive: literal tables, deterministic builders, independent
symbolic checkers, census, search code, receipts and source notices.
\end{titlepage}
\hypersetup{pageanchor=true}
\pagenumbering{roman}
\begingroup
\footnotesize
\setlength{\parskip}{0pt}
\tableofcontents
\endgroup
\clearpage
\pagenumbering{arabic}

\section{Result and exact scope}

The investigation produces five new 194-state binary blank-input Turing
machines. The principal fallback result is \textbf{PA194LC\_AX11H}, which
preserves every resolved source instruction of PA194AX11H while changing
three native transition rows. The revised AX11H packet supplies a written
absolute iPRA calibration; the new relative simulation proved here
transfers that theorem to the new table. Its full construction, source
identity, universal physical contracts, complete startup and structural
novelty have been checked.

The four other new machines are
\textbf{PA194LC}, \textbf{PA194LC\_A}, \textbf{PA194LC\_B}, and
\textbf{PA194LC\_C}. Each has exactly 194 working controls. They implement
exactly the same resolved register-machine program as the pinned primary
PA194: every one of its 3,673 primitive instructions, including every
branch and jump target, is retained. Their physical source trees and
transition graphs change. In particular, PA194LC changes a recurring
loop's binary association and counted continuation, and can be given as a
three-row patch of the primary native table.

An independently reproduced construction, \textbf{PA194GAX11}, combines an explicit
generalization binder with an eleven-entry, four-field AX command. It
changes the command protocol and proof-word compiler. Its ordinary
source/command proof and full physical checks are supplied, but the
primitive-recursive witness internalization for its relative iPRA claim
is not completed in that reproduction. During the final source check,
a newly delivered AX11G packet was found to contain exactly the same
table. GAX11 is consequently retained as a checked concurrent
reproduction, and is excluded from the five-new-graph claim.

No verified sub-194 construction emerged. The archive preserves the
specified search families and rejects candidates that achieve a numerical
state count while losing either the intended source program, a universal
physical contract, or blank initialization. The negative results are
bounded and representation-dependent; none proves a global lower bound.

\subsection{The imported absolute calibration and its proof status}

The authoritative PA README at the pinned revision states that neither
machine-specific direction of the desired iPRA equivalence is established
for the primary PA194. It distinguishes ordinary PA arguments, standard
finite-run Lean results, Python certificates, and conditional iPRA
endpoints. No table has its complete internal calibration premise
discharged in that formal development.\cite{pa-readme}

The later snapshot \texttt{603c58ca3bd6} changes the written-mathematics
situation: its revised AX11H archive supplies explicit completion papers
for the physical/source bridge, source/command bridge, current Foundation
adapter, term suppliers, and arbitrary internal PA-proof compiler, then
applies the existing one-way weak-arithmetic transfer theorems.
Its stated result is the absolute iPRA equivalence for AX11H and
QAX11H.\cite{latest-ax11h} A focused review found no identified mathematical
gap in those supplied bridge arguments; the upstream quick replay also
passed its finite bindings. These facts do not constitute a new kernel
verification of the large arithmetic proof. The original, exact upstream
archive and the audit memo are included for scrutiny.

The principal new result imports that written AX11H theorem and composes
it with the independently checked relative iPRA transfer for
PA194LC\_AX11H. Thus the requested 194-state alternative is supplied at
the ordinary written-proof level. Its exact absolute corollary and the
identity of the imported consistency predicate are stated below.

For the separate original-primary construction PA194LC, the result
established directly here is the relative theorem
\[
\ipra\vdash\Halt(M_0)\leftrightarrow\Halt(\ML),
\qquad
\ipra\vdash\neg\Halt(M_0)\leftrightarrow\neg\Halt(\ML),
\]
where $M_0$ is the exact original primary table identified below. If the independent
premise
\[
\ipra\vdash\neg\Halt(M_0)\leftrightarrow\Con(\pa)
\tag{Calibration}
\]
is established, intuitionistic composition immediately gives the requested
absolute equivalence for PA194LC and its three siblings. This is separate
from the direct import of the newly completed written AX11H calibration;
no primary-specific formal Lean endpoint is inferred merely from that
newer manuscript.

The relative proof is an ordinary mathematical proof supported by complete
finite symbolic calculations whose rules have an iPRA soundness argument.
It is not a proof checked by a newly written Lean kernel client. This is
the precise sense in which ``proved in iPRA'' is used for the relative
result in this report.

\begin{table}[ht]
\centering\small
\begin{tabularx}{\textwidth}{@{}p{.29\textwidth}XXX@{}}
\toprule
Obligation & Original-primary LC family & LC\_AX11H & Imported AX11H calibration\\
\midrule
Literal complete table & Executed & Executed & Parent pinned\\
Universal physical source contracts & Executed & Executed & Parent freshly replayed\\
Complete blank startup & Executed & Executed & Parent freshly replayed\\
Logical/source preservation & Exact resolved graph identity & Exact resolved graph identity & Written PA compiler supplied\\
Relative iPRA argument & Explicit PR witnesses and polynomial clock & Explicit PR witnesses and polynomial clock & Absolute theorem imported\\
New compiled Lean endpoint & No & No & No\\
\bottomrule
\end{tabularx}
\caption{Evidence levels. ``Executed'' refers to the included finite checkers, not to a claim that Python execution is a proof-assistant kernel.}
\end{table}

\section{Machine convention and inherited architecture}

Every released machine is deterministic, has one two-way infinite tape,
uses alphabet $\{0,1\}$, and starts with all cells zero and the head at the
origin. Each transition writes one bit, moves left or right by one cell,
and enters a new control. There are no stay-put moves. The distinguished
terminal control has no transition rows and is excluded from the state
count. In the dense tables the initial control is 0, the working controls
are $0,\ldots,193$, and the terminal label is 194. Each table therefore
has exactly 388 rows.

The native numbering used for local reasoning starts at 74, returns to
the source dispatcher at 84, and uses terminal label 254. Native labels
are identifiers, not a consecutive state count. Each directory supplies
\path{labels.json}, whose first entry is 74 and last entry is 254, and a
checked bijection between native and dense presentations.

The primary dense file has SHA-256
\begin{quote}\small
\path{4bfea62b17aad7d8c8c5e682fe8942e0a86b93e4351f8e675e32705bdaae007f}.
\end{quote}
Its native file has SHA-256
\begin{quote}\small
\path{e0a7397cf2134e8e7f58a8b1de2f7ff90b74ededd54c912f6ba5b7db48dabe54}.
\end{quote}
The frozen primary files in \path{reference/primary/} and their input
manifest make reconstruction independent of future repository changes.

The construction has four conceptually separate layers. A numerical
calculus supplies theorem codes and inference rules. A command scheduler
enumerates encoded words that manipulate a theorem, a stack, an offset,
and a cache. A nine-register source program implements the commands using
increments, conditional decrements, jumps and halt. A binary source tree
and unary-register tape engine implement those primitives. The cost is
the complete literal machine, including initialization, not the number
of logical axioms or source-tree nodes alone.

For example, the inherited coefficient-four constructors use Cantor
pairing $p(a,b)=(a+b)(a+b+1)/2+b$ and distinguish numerical forms by
residues. The primary implication and universal codes are
$I(A,B)=4p(A,B)$ and $U(a,B)=4p(B,a)+2$. A compact multiplication-by-four
source loop occurs 61 times. The new LC family changes that repeated
loop's physical decoding without changing any of those numerical
conventions, guards, proof rules, or source-register effects.

Logical registers $0,\ldots,8$ occupy physical positions
\[
(6,2,0,4,3,7,8,5,1).
\]
A value $v$ is represented by the unary block $B(v)=0\,1^{v+1}$. The
source frame also contains a fixed header allowance of 48 cells, a
nonnegative gap, and a nonempty finite spare list. The explicit gap and
spares are part of the invariant; ignoring them would invalidate a
whole-tape simulation argument.

\section{The loop-continuation constructions}
\label{sec:construction}

\subsection{The local source transformation}

Write $i$ for \texttt{inc 1}, and $c_j$ for
\texttt{continueTo $j$}. A pair $(u,v)$ is a binary source-tree node,
with left and right children $u,v$. In the primary machine the repeated
loop body, named \texttt{node62}, is
\[
T_0=(((i,i),(i,i)),c_2).
\]
PA194LC replaces it by
\[
\boxed{T_{\mathrm{LC}}=((i,i),((i,i),c_3)).}
\]
The loop still decrements register 7 and, for a positive branch, executes
four increments of register 1 before returning to the same decrement
test. One successful iteration changes $(r_1,r_7)$ to
$(r_1+4,r_7-1)$ and leaves all other logical registers unchanged.

The counted return is essential. A continuation is resolved by stripping
the specified number of right branches from the source address. Moving
the continuation one right level deeper requires changing $c_2$ to $c_3$.
Retaining its old count can produce a superficially small decoder for a
different, incorrect loop.

\begin{center}
\begin{tabular}{@{}lll@{}}
\toprule
Leaf & Primary local address & PA194LC local address\\
\midrule
First increment & \texttt{000} & \texttt{00}\\
Second increment & \texttt{001} & \texttt{01}\\
Third increment & \texttt{010} & \texttt{100}\\
Fourth increment & \texttt{011} & \texttt{101}\\
Continuation & \texttt{1}: $c_2$ & \texttt{11}: $c_3$\\
\bottomrule
\end{tabular}
\end{center}

This changes 305 primitive source addresses across the 61 occurrences.
Two independent address resolvers, based respectively on arithmetic
addresses and source intervals, recover exactly the primary ordered
instruction list, including every successor. The equivalence is stronger
than agreement of the sequence of opcodes: conditional targets and
counted return destinations are compared too.

\subsection{The exact three-row native patch}

After identifying the new internal reader with existing native label
190, PA194LC changes precisely these three rows of the folded primary
table:
\begin{center}
\begin{tabular}{@{}rrll@{}}
\toprule
Control & Read & Primary action & PA194LC action\\
\midrule
139 & 0 & $0\,R\,190$ & $0\,R\,99$\\
139 & 1 & $0\,L\,16$ & $1\,R\,190$\\
190 & 1 & $1\,R\,99$ & $1\,L\,73$\\
\bottomrule
\end{tabular}
\end{center}
The row $(190,0)=0\,R\,99$ is retained. Every other native row is
identical. These labels refer specifically to the pinned primary native
presentation. Applying the same integers to an arbitrary dense table
would be a different edit.

The raw construction has 201 working controls, and the same seven
simultaneous contextual folds yield 194. Each fold identifies two reader
roles whose incoming transitions leave opposite bits immediately behind
the head. The merged control recovers that history locally. All four
combinations of history and current bit are checked for each fold in
the final simultaneously folded table, including equality of the
modified tape cells, endpoint control and head displacement.

The patch rows participate in 122, 183 and 61 symbolic source outcomes,
respectively. They change the implementation of recurring computations;
they are not unused added states. This statement concerns universal source
contracts. It does not assert that every changed row has appeared in a
short blank run. In fact none is used in the fixed startup prefix.

\subsection{Three further source-preserving graphs}

The other released trees are
\begin{align*}
T_A&=(i,((i,(i,i)),c_3)),\\
T_B&=((i,(i,i)),(i,c_3)),\\
T_C&=(((i,i),i),(i,c_3)).
\end{align*}
Their raw tables have 202 controls and admit eight checked contextual
folds. All retain the same resolved source graph and the same nine logical
registers. They differ from one another and from PA194LC even after a
strong structural comparison that permits arbitrary state renaming,
any spanning initial root, global left/right reflection, global bit
exchange, and ignored write/move decoration on a final halting edge.
The final comparison finds no match among the 110 prior complete
194-state graphs. Symbol exchange here strengthens a novelty test; it is
not asserted to preserve the chosen zero-blank starting convention.

\begin{table}[ht]
\centering\small
\begin{tabular}{@{}lrrrrr@{}}
\toprule
Machine & Raw & Folds & Final & Changed addresses & Startup\\
\midrule
Primary & 201 & 7 & 194 & 0 & 58,154\\
PA194LC & 201 & 7 & 194 & 305 & 58,154\\
PA194LC\_A & 202 & 8 & 194 & 305 & 58,320\\
PA194LC\_B & 202 & 8 & 194 & 305 & 58,320\\
PA194LC\_C & 202 & 8 & 194 & 183 & 58,320\\
\bottomrule
\end{tabular}
\caption{Exact reconstruction counts. Startup is counted in physical TM transitions. Maximum primitive address length is 36 in each new construction.}
\end{table}

These are alternative implementations, not speed improvements. With
matched input frames, a complete successful multiplication iteration in
PA194LC costs six more physical transitions than the primary. The
symbolic checker's operation count is a verification workload and must
not be interpreted as a running-time comparison of the Turing machines.

\subsection{Independent physical verification}

The builder checks the inherited input hashes, compiles the selected tree,
applies the recorded folds, creates the dense numbering, compares the
source graph, and simulates startup. A separate checker imports neither
the builder nor its source-flattening routine. It reads the literal
native and dense files, checks their conjugacy, resolves the source in two
independent ways, derives four list-traversal lemmas from the literal
transition rows, and proves every primitive's complete frame contract.

Each positive decrement is represented by a successor-valued register;
the zero branch is checked separately. The parameters range over nine
arbitrary natural primary registers, an affine gap and an arbitrary
nonempty finite spare list. The interpreter accelerates only justified
constant-run and generic-list traversals. It uses no opaque arithmetic
box oracle. All 4,780 outcomes close at the required exact endpoint,
with strictly positive duration. Complete blank-tape startup equality is
checked independently of those arbitrary-frame contracts.

\begin{center}
\begin{tabular}{@{}lrrr@{}}
\toprule
Table & Source outcomes & Symbolic operations & Largest case\\
\midrule
Primary, fresh replay & 4,780 & 1,078,198 & 397\\
PA194LC, final native labels & 4,780 & 1,077,405 & 397\\
PA194LC\_A & 4,780 & 1,083,389 & 397\\
PA194LC\_B & 4,780 & 1,083,023 & 397\\
PA194LC\_C & 4,780 & 1,082,636 & 397\\
\bottomrule
\end{tabular}
\end{center}

The archived receipts bind each table's dense and native SHA-256 to its
contract file. The reported counts are observed completed runs. The
code's soundness argument and the arithmetization of composition are
given in the next section. A fresh invocation with
\texttt{python3 reproduce.py --physical} regenerates the contracts as well
as checking the finite reconstruction; the faster default mode verifies
the recorded certificate bindings and reruns the other finite checks.

\section{An explicit quadratic relative simulation in iPRA}
\label{sec:quadratic-transfer}

Let $M_0$ denote the pinned primary PA194 table and let $M_L$ denote
PA194LC.  This section proves a relative result.  No assertion identifying
either machine's halting with a PA refutation is used as a premise.
For a literal table $M$, totalize its transition function by leaving a
halting configuration fixed, and write
\[
 H_M(t)\;:\!\iff\;
 \text{the configuration after $t$ transitions from the blank start is halted}.
\]
This is a decidable primitive-recursive predicate; in the totalized
semantics it means that the machine has halted by time $t$.

The exact startup and local-contract certificates for the two tables,
together with the symbolic-rule soundness argument below, give the
following ordinary computer-assisted arithmetic theorem:
\[
 \mathsf{iPRA}\vdash\forall t\,
 \bigl(H_{M_0}(t)\longrightarrow H_{M_L}(\Phi(t))\bigr),
 \qquad
 \mathsf{iPRA}\vdash\forall t\,
 \bigl(H_{M_L}(t)\longrightarrow H_{M_0}(\Phi(t))\bigr),
 \tag{Q}\label{eq:quadratic-transfer}
\]
where the explicit polynomial is
\[
 \boxed{\Phi(t)=58\,154+13(t+1)^2+681(t+1)
              =13t^2+707t+58\,848.}
\]
Consequently iPRA proves equivalence of their halting sentences, and
intuitionistic logic gives equivalence of their nonhalting sentences.
This is not a newly compiled Lean theorem.  The proof's mathematical
hypotheses are the complete literal-table contracts and their stated
arithmetical soundness, not the unfinished absolute consistency
calibration of the original primary table.

\subsection{Primitive-recursive configurations and local proofs}

Use the canonical tape zipper $(q,l,r)$, with the low bit of $r$ at the
head and the low bit of $l$ immediately to its left.  A write of $w$ followed
by a right move sends $(l,r)$ to
\[
 (2l+w,\lfloor r/2\rfloor),
\]
and a write followed by a left move sends it to
\[
 (\lfloor l/2\rfloor,
   4\lfloor r/2\rfloor+2w+(l\bmod 2)).
\]
This represents the entire finite-support tape, including arbitrary blank
exteriors, by natural numbers.  Table lookup, the totalized step, finite
word operations, and bounded iteration are primitive recursive.  Their
composition equations require only quantifier-free induction.

For each implementation augment the logical source configuration by a
nonnegative gap $g$, a first spare value $a_0$, and a finite spare tail
$A$.  The frame function $E(s)$ is thereby a primitive-recursive function,
not a relation requiring an unbounded existential witness.  The source
counter selects one of the finite list of valid instruction addresses;
the nine logical registers and the auxiliary fields are explicit coded
parameters.  Every frame has the nonhalting root control.  In particular,
a frame whose next source instruction is \emph{halt} is still a
nonhalting physical configuration before that instruction is executed.

The independent checker derives four generic list traversals directly
from the literal rows.  Their proofs use only: the defining equation for a
single transition; quantifier-free induction for a self-loop crossing a
constant-symbol run; finite coded-word identities; and induction on a
coded list's length.  If
\[
 B(v)=0\,1^{v+1},\qquad C(v)=1\,0\,1^v,
\]
the normalization identity $C(A)1=1B(A)$ follows from the corresponding
single-block identity by list-length induction.  The list traversal has
exact duration $\operatorname{width}(A)=\sum_{v\in A}(v+2)$.
Exterior-word codes remain free parameters throughout these proofs.
Canonical zipper equality replaces an unbounded assertion about every
tape coordinate.

The finite symbolic proof for each source site composes these lemmas.
Every decrement has its separate zero and successor-value branches.
All intermediate run lengths have coefficientwise nonnegative affine
expressions, and every complete contract has strictly positive duration.
The checker does not use arithmetic-box oracles.  Its finite successful
execution is evidence for the explicit derivations generated by these
sound rules; Python is not being treated as a proof-assistant kernel.

The two independent address resolvers agree with each other and with all
$3\,673$ entries, including jump targets, of the primary resolved source
program.  The implementations can have different physical paths, gaps,
and auxiliary updates.  Their logical projections are nevertheless the
same primitive-recursive source iteration, by quantifier-free induction.

\subsection{Synchronization and bounded extraction}

For one implementation, let $h(s)$ be the Boolean saying that the source
counter is at its explicit halt instruction.  Let $S$ be the augmented
source step, with a common absorbing logical halt extension and arbitrary
primitive-recursive auxiliary updates after such an instruction.  Write $s_n=S^n(s_0)$.  Let $k(s)\ge1$ be the exact
normal-instruction duration, and $k_H(s)\ge1$ the halt-instruction duration;
extend both functions totally on unused inputs.  If $b$ is the fixed
startup time, define
\[
 T(0)=b,\quad T(n+1)=T(n)+k(s_n),\qquad
 B(n)=\prod_{j<n}(1-h(s_j)).
\]
The Boolean $B(n)$ and clock $T(n)$ are primitive recursive.
The synchronization invariant is the quantifier-free implication
\[
 B(n)=1\ \longrightarrow\
 \bigl(\operatorname{Inv}(s_n)\ \wedge\
   \operatorname{run}_M(T(n),\mathrm{blank})=E(s_n)
   \ \wedge\ T(n)\ge b+n\bigr).
 \tag{S}\label{eq:source-synchronization}
\]
The startup is its base case.  The defining equation
$B(n+1)=B(n)(1-h(s_n))$ supplies both the predecessor invariant and the
nonhalt-source premise required by the next local contract.  Exact
step-iteration composition gives the physical equation.  Positivity of
$k$ gives the clock inequality.  The frame invariant is preserved because
the source targets explicitly preserve a first spare, keep the gap
nonnegative, and select another valid source address.  Thus the induction
formula contains no unbounded quantifier over runs, tapes, or proofs.

For the forward direction, a source witness $h(s_n)=1$ can be replaced by
the least $m\le n$ with that property, using bounded minimization.  Then
$B(m)=1$, and the halt-site contract gives a physical halt at
$T(m)+k_H(s_m)$.  This first-occurrence step avoids silently treating an
arbitrary source-halting witness as the first one.

For the reverse direction suppose $H_M(t)$.  If $B(t+1)=1$, equation
\eqref{eq:source-synchronization} gives a nonhalting root frame at time
\[
 T(t+1)\ge b+t+1>t.
\]
Absorbing halting implies that the machine is halted at this later time,
a contradiction.  Therefore $B(t+1)=0$.  A bounded first-hit function
returns some $m<t+1$ for which $h(s_m)=1$.  Its correctness is proved by
quantifier-free induction on the finite Boolean prefix.  In particular,
\[
 m\le t.
\]
There is no unbounded search, inverse decoding of intermediate tapes,
Markov principle, or elimination of double negation for an unbounded
existential.  The indexing convention matters: $h$ marks a
\emph{ready-to-execute} halt instruction, whose physical frame is still
nonhalting.  Under a different convention one may use the prefix through
index $t+1$ and the harmless bound $m\le t+1$ instead.

A hidden halt inside a claimed normal segment cannot defeat the
argument: in absorbing semantics a correct equation ending at a
nonhalting frame already excludes such an earlier halt.

\subsection{A uniform polynomial clock}

The actual complete startup tapes of both tables were separately
simulated and decoded.  After $b=58\,154$ transitions each is at native
control $84$, head $-2$, with source address $100000$ (instruction $25$),
and with
\[
 g=0,\qquad (r_0,\ldots,r_8)=(2,0,\ldots,0),
 \qquad a_0=1,\qquad A=[1,\ldots,1]\text{ of length }25.
\]
There are $26$ spare blocks altogether, so
$\operatorname{width}(A)=25(1+2)=75$.
The replay compares the entire finite support, not merely the control,
head position, or a bounded tape prefix.

For each of the $4\,780$ source outcomes of each literal table, a finite
coefficientwise comparison of its thirteen-slot duration vector gives
\[
 3\le k(s),k_H(s)\le
 534+8g+4\sum_{i=0}^{8}r_i+2a_0+2\operatorname{width}(A).
 \tag{D}\label{eq:affine-duration-envelope}
\]
Only the applicable normal or halt duration is used in a given case.
A successor-decrement certificate uses $r_i-1$ for its variable slot,
which can only decrease this upper bound.  The machine and contract
hashes are bound to the coefficient receipt.

Each source primitive increments at most one logical register by one.
Its gap update is $g'=g+\delta+o$, where both $\delta$ and the address
overflow bit $o$ belong to $\{0,1\}$.  Increment appends one spare value
$1$; every other primitive retains the spare list.  Consequently,
after $j$ normal instructions,
\[
 \sum_i r_i\le2+j,\qquad g\le2j,\qquad
 a_0=1,\qquad \operatorname{width}(A)\le75+3j.
\]
These are again quantifier-free induction invariants of the explicit
primitive target-frame updates.  Substitution in
\eqref{eq:affine-duration-envelope} gives
\[
 k(s_j),k_H(s_j)\le
 534+16j+4(2+j)+2+2(75+3j)=694+26j.
\]
Hence $n$ normal instructions take at most
\[
 b+\sum_{j<n}(694+26j)=b+13n^2+681n,
\]
and a first halt instruction at index $n$ yields halting by
\[
 b+13(n+1)^2+681(n+1).
\]
Given a physical halt at time $t$ in either implementation, bounded
extraction supplies the same logical halt-site index $n\le t$ for the
other.  Its forward bound and absorbing halting prove
\eqref{eq:quadratic-transfer}.  The affine upper bound is stronger than
needed for relative iPRA provability---total primitive-recursive clocks
and positivity already suffice---but makes both witness transformations
fully explicit.

The proof contains no compiler from arbitrary PA proofs to numerical
refutations.  That separate calibration is still needed to replace one
of the halting sentences by PA inconsistency.  Once it is established
for the primary table in iPRA, the relative theorem transfers it to
PA194LC immediately.


\subsection{Transfer to the three siblings}

The same coefficient envelope and frame-growth inequalities hold for
PA194LC\_A, PA194LC\_B and PA194LC\_C. Their whole-tape startup frames
have the same source ordinal, register values, gap, and spare list, but
take 58,320 transitions. Thus for any ordered pair $M,N$ among the
primary and the four new source-preserving tables,
\[
\ipra\vdash\forall t\,
\bigl(H_M(t)\longrightarrow H_N(\Psi(t))\bigr),
\qquad
\boxed{\Psi(t)=13t^2+707t+59\,014.}
\]
The larger constant is a uniform convenience; for a target with the
shorter startup, $\Phi$ suffices. No multiplication of the two clocks is
needed: bounded extraction first finds the common source halt-site
index, then the target's forward clock supplies the witness.

\section{PA194LC\_AX11H and the requested consistency equivalence}
\label{sec:lc-ax11h}

The final construction applies the same loop-continuation transformation
to the revised AX11H packet's primary machine. Denote that parent by
$M_H$ and the new machine by $\MLH$. The parent dense table is bound by
\begin{quote}\small
\path{c5dfe1026d077d3e062eca11edb789597bfbbced4f6040f8828ae5d467b1167a},
\end{quote}
and the new dense table by
\begin{quote}\small
\path{5f10f071e7118ff33aba822e05bae0331b4c78508656c3e945c31b0c6788f9c5}.
\end{quote}
The full new table is printed in Appendix~\ref{app:lc-ax11h-table} and
provided in \path{additional/PA194LC_AX11H/machine.tm}.

\subsection{Exact construction and completed checks}

The parent uses the four-field AX11 command, an eleven-entry saved prefix,
visible no-push MP, and the H halt test: Shift halts precisely when the
cached theorem is code one. The loop transformation leaves that command
protocol and its numerical calculus unchanged. All 3,671 resolved
instructions, including every register operand and successor, are
identical. The 61 recurring constructor loops change association and
replace their depth-two continuation by a depth-three continuation;
305 primitive addresses change.

Both the raw and final folded native tables differ in exactly the three
rows displayed in Section~\ref{sec:construction}: $(139,0)$ becomes
$0R99$, $(139,1)$ becomes $1R190$, and $(190,1)$ becomes $1L73$.
Here the parent of the patch is $M_H$, not the original primary $M_0$.
All seven contextual folds remain valid, giving $134+67-7=194$ working
controls and 388 transition rows. Initial and terminal controls are 0
and 194 in the dense numbering. The complete blank startup takes 57,922
transitions in both implementations.

\begin{center}
\begin{tabular}{@{}lrr@{}}
\toprule
Independent check & Parent AX11H & New LC\_AX11H\\
\midrule
Resolved source instructions & 3,671 & 3,671\\
Universal branch outcomes & 4,778 & 4,778\\
Symbolic operations & 1,101,836 & 1,101,043\\
Maximum operations in one case & 403 & 403\\
Minimum duration constant & 3 & 3\\
List lemmas derived from literal rows & 4 & 4\\
Whole-tape startup transitions & 57,922 & 57,922\\
\bottomrule
\end{tabular}
\end{center}

These are fresh full literal-row replays of both tables. The separate
construction checker also passed all 4,778 outcomes, using 1,514,306
symbolic operations, maximum 506. Both independent address resolvers
agree; complete source identity, native/dense conjugacy, simultaneous
fold cases, and full startup tape equality pass. The changed successful
loop iteration costs exactly six additional physical transitions at each
of its 61 occurrences. All five newly delivered LC graphs are pairwise
distinct and match none of the 110 prior graphs in the latest audited
snapshot. GAX11's separate duplicate is deliberately excluded.

\subsection{An explicit relative iPRA theorem}

Let $H_M(t)$ again mean halted by time $t$ in absorbing terminal semantics.
The finite coefficient audit of all 4,778 contracts in each table gives
the same envelope as in Section~\ref{sec:quadratic-transfer}:
\[
 3\le k\le534+8g+4\sum_{i<9}r_i+2a_0+2\operatorname{width}(A).
\]
The actual maximum vector is
$(534,8,4,4,4,4,4,4,2,4,4,2,2)$.
The exact startup frames have $g=0$, logical registers $(2,0,\ldots,0)$,
first spare $a_0=1$ and 25 further unit spares. Consequently the same
primitive updates give the same bound $k(s_j)\le694+26j$.

\begin{theorem}[Relative AX11H simulation]
With
\[
 \boxed{\Theta(t)=57\,922+13(t+1)^2+681(t+1)
                    =13t^2+707t+58\,616,}
\]
the following implications are provable in iPRA:
\[
 \forall t\,[H_{M_H}(t)\to H_{\MLH}(\Theta(t))],\qquad
 \forall t\,[H_{\MLH}(t)\to H_{M_H}(\Theta(t))].
\]
In particular iPRA proves equivalence of the two blank nonhalting
sentences.
\end{theorem}

\begin{proof}
Use the canonical zipper, explicit augmented source iteration, common
absorbing logical halt extension, Boolean prefix $B(n)$ and exact
clock $T(n)$ from Section~\ref{sec:quadratic-transfer}, now instantiated
with the 3,671-instruction AX11H source. Each local contract is proved by
the same literal-transition, unary-run and finite-list rules. The
independent checks establish all source targets and exact frame updates
for these two specific tables, so the quantifier-free synchronization
invariant applies with startup constant 57,922.

A physical halt by $t$ gives, by the bounded first-hit argument, a
common source halt-site index $n\le t$. In the other implementation,
the first $n$ normal instructions and its final halt instruction take
at most
\[
57\,922+\sum_{j=0}^{n}(694+26j)
=57\,922+13(n+1)^2+681(n+1).
\]
Monotonicity and absorbing halting yield the displayed witness at
$\Theta(t)$. All constructions and bounded Boolean invariants are
primitive recursive. No unbounded witness search or Markov principle
is used. The argument in both directions is identical after exchanging
the tables. A self-contained expanded version is included in the new
machine's standalone \path{proof.tex} and PDF.
\end{proof}

\subsection{Importing the revised absolute calibration}

The original upstream archive at the final snapshot is preserved without
modification as
\path{reference/ax11h_calibration_603c58c.zip}; its SHA-256 is
\begin{quote}\small
\path{93ee082ef035b2d32e71d6af5239b25b4bea1773f03bea4d2c39810a7fe15075}.
\end{quote}
Its \path{proof/consistency_equivalence.tex} states and proves, using its
separate completion papers and the pinned generic transfer results,
\[
 \ipra\vdash\neg\Halt(M_H)\leftrightarrow\Con(\pa).
 \tag{AX11H calibration}
\]
Here $\Con(\pa)$ is precisely the packet's
\texttt{nativePAConsistency}. This negates \texttt{nativePARefutation}
for Foundation dependency
\path{32e1a0956a8622fad067328ca1959729a7634428}.
Foundation falsity code 13 and numerical-calculus contradiction code 1
are connected by explicit proof translations; they are not identified.
\cite{latest-ax11h}

The imported argument supplies more than an assertion that a standard
simulation can be internalized. Its Foundation adapter gives explicit
proof opening/closing constructors, bounded induction-axiom witnesses,
and the correction for the changed finite order axiom. Its term and
formula compilers use acyclic primitive-recursive supplier strata and
bounded demand trees, with Boolean correctness invariants for the
actual requests, instead of full-PA induction over all environments.
Its source-command paper binds the literal arithmetic templates and
four-field parser, and its physical paper proves frame closure and
the symbolic rules uniformly. The final paper composes the two
all-model $I\Sigma_1$ directions with the existing one-way transfer
endpoints. The packet does not implement the entire large PA proof
compiler as executable code and does not supply a new end-to-end
Lean or Rocq build.

The focused review retained in \path{verification/ax11h_absolute_audit.txt}
found no identified gap in those completion arguments. The upstream
quick verifier passed all 1,186 manifest files, byte-identical
reconstruction of both parent designs, source-command and physical-frame
audits, fixed proof constructors, 119 current-Foundation coding
comparisons, and the premise-exchange certificate. That quick mode
explicitly does not verify the entire arithmetic derivation; the full
literal physical replays for $M_H$ and $\MLH$ were performed separately
as reported above.

\begin{theorem}[The 194-state fallback, importing the AX11H theorem]
Using the revised packet's written AX11H calibration as an imported
mathematical theorem,
\[
 \boxed{\ipra\vdash\neg\Halt(\MLH)\leftrightarrow\Con(\pa).}
\]
\end{theorem}
\begin{proof}
The relative theorem gives
$\neg\Halt(\MLH)\leftrightarrow\neg\Halt(M_H)$ in iPRA.
Compose it intuitionistically with the imported AX11H calibration.
\end{proof}

This explicitly identifies the inherited mathematical dependency.
The new contribution is the different literal machine and its exact,
positive-duration iPRA transfer, with complete source and tape evidence.
The corollary is a written proof using that supplied calibration, not
a claim that the unfinished primary-machine Lean development has acquired
a new kernel-checked endpoint.


\section{Comparison of the supplied designs}
\label{sec:corpus-comparison}

The comparison is pinned to ProveIt commit
\texttt{32d0d8b7d93a63ad4106c18cdb29b37927acdad2}.
It covers the 77 filed research packets under the PA directory and the five
ZIP deliveries incoming at that initial snapshot in \nolinkurl{docs/incoming}. The repository's
primary documentation distinguishes experimental tables, written
arithmetical arguments, executable certificates and completed formal
theorems; the same distinctions are retained here. In particular, the
state count of a table does not certify its consistency property.
\cite{pa-readme,machines}

\subsection{Census method and what the counts mean}

The fresh census reads the pinned Git blobs, verifies their Git SHA-1
identities and recursively inspects ZIP members. It recognizes complete
numeric five-column transition tables and the documented explicit JSON
formats: row arrays, paired actions, row records, transition dictionaries
and embedded whole-table strings. It does not run arbitrary generators
or reconstruct a table from prose, TeX or Lean definitions. This is an
explicit, reproducible recognizer scope rather than a claim to have
decided what every file could compute.

\begin{center}
\begin{tabular}{@{}lr@{}}
\toprule
Freshly counted object & Count\\
\midrule
Tracked source inputs verified & 12\,551\\
Direct and archive-member payloads inspected & 15\,172\\
Recognized complete table occurrences & 3\,606\\
Normalized labeled tables & 778\\
Complete spanning one-halt graph classes & 658\\
Such graph classes with 194 working states & 94\\
Such graph classes with 193 working states & 53\\
\bottomrule
\end{tabular}
\end{center}

A counted graph has two left/right transitions per working control,
exactly one external halt target, and some working root reaching its
entire action graph. Six normalized nonspanning tables are retained
separately; 50 incidental one-state JSON fragments are excluded. The
archive traversal sees 19 distinct ZIP payloads through 31 references.
No source-identity or parse errors were reported. Full receipts and
recognizers accompany the report in \nolinkurl{comparison/}.

The official filed-table index uses a different scope: its 662 classes
include some nonspanning components and partial engines, whereas the
present census includes explicit JSON and the initial incoming ZIPs. Its 89
classes at 194 reachable states therefore cannot be added to the
incoming reference collections. The fresh census finds five additional
complete 194-state graphs in the initial incoming packets, giving 94.
\cite{machines}

Canonicalization removes working-control names. The independent novelty
comparison also permits any spanning root, left/right reflection,
global symbol exchange and arbitrary final-edge write/move decorations.
It compares entire canonical tuples, not only their hashes. All 94
supplied 194-state classes remain distinct under that larger equivalence.
Symbol exchange makes the structural comparison more demanding; it is
not used to claim preservation of a specified zero-blank start.
Likewise, graph reachability is not coverage of the actual blank run.
Some native HAZ files have dense labels but a designated start other
than zero; start metadata remains separate from graph identity.

Appendix~\ref{app:frontier194} identifies every one of the 94 graphs;
Appendix~\ref{app:packet-ledger} identifies every delivery. Their attributed
categories are 51 principal constructions, seven supplementary
constructions, 25 research auxiliaries, four search seeds and seven
rejected constructions. These are evidence categories, not 51 completed
iPRA consistency theorems. Packet mechanisms and historical proof
claims below are attributed to that ledger unless a fresh check is
expressly identified.

\subsection{The architecture and the principal sources of variation}

The common construction has four interacting levels: the numerical
proof calculus; the register program that enumerates command words and
manages the store; the source tree and its binary readers; and the unary
tape engine with blank initialization. Changes at different levels have
different proof obligations. A shorter formula program need not produce
a smaller physical table, and two advertised savings may consume the
same decoder equality.

\paragraph{Baseline and physical placement.}
Primary PA194 has a separate diagonal emitter, a ten-entry AX stack
prefix and the original sum-testing Shift guard. Its physical allocation
is $134+67-7=194$: fixed controls, raw AX readers, then contextual folds.
The alternate decoder has $134+68-8=194$. AD exchanges the physical A/D
register positions. TDY cycles three positions, P5 moves five registers
and seven router targets, and HX exchanges H/X. These preserve the
abstract resolved source but change its physical realization. They
require complete tape-frame and initialization checks, even when their
logical proof compiler is unchanged.\cite{pa-readme}

\paragraph{The store and the contradiction guard.}
C caches the result of successful modus ponens and makes Shift test that
cached head; CS uses a symmetric destructive comparison. N instead
removes the dummy-zero push, leaving the consequent visible, and keeps
the original sum guard. Its proof-word compiler needs a reserved zero.
H combines that visible-result protocol with a cache-only guard; HS
uses the symmetric version. SUMS and the uppercase NS design retain
the original sum test and make its comparison symmetric. This distinction
matters: the lowercase \nolinkurl{pa194ns} packet contains the HS graph,
whereas uppercase NS denotes a different, sum-guard graph.

\paragraph{Logical heads and formula arithmetic.}
Q exchanges the two curried premises of the universal-distribution
head. QN, QNS, QH and QHS combine this change with the corresponding
store and guard. The logical transfer needs a primitive derivation of
premise exchange. G changes the saved generalized head from $U(d,h)$ to
$U(c,h)$, binding the old cached theorem with AX's explicit first
operand; the generalization compiler changes accordingly. D emits
$\Gamma(D)$ in place of $\Gamma(A)$; DHS combines that change with the
no-push, symmetric cached controller, at $135+69-10=194$ controls.

The coefficient-five and coefficient-six families---M5/M6, N5/N6,
H5/H6, Q5/Q6, QN5/QN6 and QH5/QH6---change the arithmetic formula
constructors throughout the source, including 61 constructor bodies in
the supplied implementations. They require a logical recoding proof:
unused residues, the zero-code alias and every emitted head must be
accounted for. They are not register permutations or control renamings.

\paragraph{Known-zero returns and source associations.}
J, F, A, CSZ/X, ZR and HAZ replace selected unconditional jumps with
decrements of registers known to be zero at those source locations.
The proof is an inductive invariant of reachable source states. Its
scope does not automatically include blank initialization, during
which the physical controls run outside that source invariant.
R/M4R and the saved source rotations change binary associations and
instruction addresses while preserving a resolved source graph.
Many saved rotations are search inputs rather than principal releases.
PA1944L associates four increments as $((i,(i,i)),i)$, adding one helper
and one fold to give $135+67-8=194$.

\paragraph{The AX11 interface.}
AX11 pushes its final guarded head and falls through into the full MP
parameter reader. Its command consumes a fourth field $z$ and creates
an eleven-entry prefix. The extra field is semantically real:
completeness may choose $z=0$, but soundness must handle every $z$ and
the Cantor zero-stack alias. AX11QH5, AX11Q6, AX11H, AX11QH and AX11S
combine this changed interface with earlier arithmetic, logical and
controller changes. Counting the extra push while overlooking its
following MP operation would give the wrong command semantics.

\paragraph{A genuinely different decoder protocol.}
B is a physical rewrite with a different proof shape: it merges two
history-sensitive reader roles and reuses the freed control as a
dedicated selector. Its four changed rows support a short-step,
all-tape transfer. This illustrates why failure of ordinary control
quotients leaves room for contextual designs. The loop-continuation
constructions in this report instead preserve all resolved source
instructions and alter the placement of the continuation within a
repeated source tree.

\paragraph{Larger ancestors and alternative schedulers.}
Original PA195 has a bundled 191-token diagonal emitter; PA195S uses
the separate 190-token interface. PA195C, PA195DC and PA195S8 make
further controller, diagonal or scanner changes. MP196 detects
contradiction immediately after a successful MP instead of waiting for
Shift. The even-candidate scheduler with H/X placement costs 195
controls. These are relevant architectural alternatives, but the supplied
tables do not meet the requested strict improvement below 194.

\subsection{Duplicates and the five initial incoming packets}

Different packet titles do not necessarily mean different designs.
The principal PA194C table occurs in five deliveries. QN equals NQ;
the two QNS deliveries agree; N6 is repeated; CSZ equals X. QH5 appears
both separately and as an AX11 companion, while QN6 appears alone and
beside QN5. The initial incoming AX11QH5 is byte-identical to its
already filed counterpart: its dense table has SHA-256
\nolinkurl{f7f52501055ff4fe47aab51acc6536adf3823afca4016a670c67b2c592ed4b02}.

\begin{center}
\small
\begin{tabularx}{\linewidth}{@{}l l X@{}}
\toprule
Packet & New graph & Distinguishing change\\
\midrule
P78 & PA1944L & Four-increment association; primary source retained.\\
P79 & PA194AX11H & AX11 with the no-push cached-head controller.\\
P80 & PA194AX11QH & AX11 with Q premise exchange and H guard.\\
P81 & PA194AX11S & AX11 with symmetric original sum guard.\\
P82 & PA194HX & Physical H/X exchange; AX11QH5 is a duplicate.\\
\bottomrule
\end{tabularx}
\end{center}

All five original deliveries explicitly distinguish their preservation arguments
from the unfinished absolute iPRA calibration. The table lists their
new graphs, not all reference tables or experiments carried inside their
ZIPs. The machine-readable census retains each occurrence and its
source locator, so duplicate evidence remains auditable without
inflating the count of designs. The later revised AX11H written proof
is treated separately in Sections~\ref{sec:lc-ax11h} and~\ref{sec:late-snapshot}.

\subsection{Why the existing 193-state objects do not settle the problem}

The 53 complete spanning 193-state graphs are not 53 accepted
consistency machines. The most direct arithmetic shortcut uses
coefficient three without a sufficient change of calculus. It identifies
two incompatible formula codes:
\[
 E_3(1,0)=I_3(0,1)=6.
\]
A supplied primitive certificate derives contradiction code 1 in
73\,393 proof lines. That number counts proof lines, not physical TM
steps. Later supplied obstructions cover all constant offsets and
cofinite cases within specified coefficient-three families; they do
not rule out every possible new formula encoding.

Runtime193 fails for a different reason. It redirects reader 301 to 88
and changes native rows
\[
 (88,0)\mapsto(0,L,51),\qquad
 (51,0)\mapsto(0,R,122).
\]
Its initialized source contracts survive, but its inherited blank start
escapes into an infinite outward zero ray. The supplied exact allocation
obstruction covers all $9!=362\,880$ register permutations under the
specified selector changes and fixed initial control; it is not a
theorem about arbitrary bootstrap redesigns. Subsequent append-zero,
transplantation, relay and scheduler variants remain rejected or lack a
complete startup calibration.

The newer even-runtime194 diagnostic is also rejected. Its supplied
certificate reaches native state 35 at head position 15 after 93
transitions, with ones only at $-3$ and 1, then follows the zero-loop
rightward forever. The new scheduler is never reached. This is a
194-state failed experiment, not an additional accepted design.

Two finite negative results help delimit the next search. First, fresh
partition refinement on all 94 graphs finds 194 singleton classes in
the greatest unconditional action-preserving control congruence of
each table. No ordinary one-step state quotient saves a control.
Second, a supplied exact rational dual bounds the net saving by seven
for fresh direct selectors and disjoint raw-action reader folds of the
fixed 201-control primary parent. Overlapping folds, new equalities
after aliasing, helper chains, longer detours and changed grammars lie
outside that theorem. Neither result is a global 194-state lower
bound. The wider search receipts in this report likewise describe
bounded experiments, not an impossibility proof.

\subsection{The arithmetic boundary and the related research programs}

At the initial pinned revision, the authoritative PA README explicitly leaves
both machine-specific directions of the absolute physical-table
iPRA equivalence open. A standard finite-run theorem, a Python
certificate and a conditional iPRA endpoint do not by themselves
discharge that calibration. Written PA equivalences in historical
packets also require a separate argument before their base theory can
be lowered to iPRA.\cite{pa-readme}

Substantial formal progress is already present. In all models with an
$I\Sigma_1$ reduct, the modeled scheduler has fairness and a
target-derivation correspondence. Guarded-gate and separation results
give the modeled scheduler's soundness direction. Generic
$I\Sigma_1$ transfer and the final intuitionistic bridge are available.
At that initial revision, the formal primary-machine development still
requires the uniform correspondence between physical
tape execution and controller/source frames at internal times, and
the compilation of arbitrary internal PA refutations into the numerical
calculus. Fixed blank startup is checked; model-length simulation is a
different obligation. No literal table's complete \texttt{Calibrated}
premise is discharged in that development. The later written AX11H
calibration is separate mathematical evidence, reviewed and imported in
Section~\ref{sec:lc-ax11h}.\cite{pa-readme}

The relative theorem for PA194LC is therefore useful on its own: it
transfers halting between two literal tables with a concrete polynomial
clock, without assuming either table's absolute calibration. Composing
it with a future completed primary calibration would establish the
requested Con(PA) equivalence. This report does not replace that
missing premise with standard-model testing. The additional G/AX11
composition has a further source-calculus transfer obligation and is
labeled separately from the instruction-preserving loop family.

The suggested adjacent projects are considered separately below, with their
size metrics and arithmetic bases kept distinct.

The fresh work supporting this comparison is the source-identity census,
canonical-tuple novelty comparison and unconditional-congruence check.
Historical startup durations, source-contract totals, proof-line counts
and packet calibrations are attributed evidence unless separately
replayed in a construction's verification section. This boundary keeps
the comprehensive comparison reproducible without suggesting that all
82 original theorem developments have been independently rebuilt.


\section{Late snapshot addendum: the comparison extended to 110 graphs}
\label{sec:late-snapshot}

During preparation, the repository advanced from the initial pinned
revision to
\nolinkurl{603c58ca3bd6bf3222fcfc4e308e6b079f65d573}.
The initial 82-delivery ledger and 94-graph census are retained with their
original provenance. This addendum separately accounts for the
\href{https://github.com/VladimirReshetnikov/ProveIt/tree/603c58ca3bd6bf3222fcfc4e308e6b079f65d573}{later snapshot}.
The five former incoming packets were filed under PA/docs. Five incoming
ZIP payloads are now present: four new archive names and a revised
AX11H delivery. The older filed AX11H and its revised incoming version
are distinct evidence records.

\subsection{The added material and the repeated census}

The same complete literal/explicit-JSON recognizers were rerun against
the later Git objects, without changing the initial census. They verified
13\,110 tracked source inputs and inspected 15\,485 direct or archive
payloads, finding 3\,128 table occurrences, 805 normalized tables and
680 complete spanning one-halt graph classes. There were no source or
parse errors. All previously recognized graph classes are still
represented. The 194-state count increases from 94 to 110; the 193-state
count remains 53. These are recognizer-scope counts, not counts of
completed consistency theorems.

Five of the 16 added 194-state graphs are reported principal constructions:
AX11G, AX11QH4L, AX11QH4R, GH and GHS. The other eleven are three B/4L
compositions used as local-search inputs and eight GH neutral-edit
search inputs. The AX11QH4 packet explicitly declines to promote its
three auxiliary B compositions to additional proved machines. The
revised AX11H packet adds no new 194-state graph relative to the
original 94: both of its principal graphs were already represented.

\begin{center}
\footnotesize
\setlength{\tabcolsep}{2pt}
\begin{tabular}{@{}>{\raggedright\arraybackslash}p{.05\linewidth}>{\raggedright\arraybackslash}p{.52\linewidth}>{\raggedright\arraybackslash}p{.23\linewidth}>{\raggedright\arraybackslash}p{.13\linewidth}@{}}
\toprule
ID & Incoming ZIP & SHA-256 prefix & Change\\
\midrule
N1 & \nolinkurl{pa194_ax11g_research_2026-11-14.zip} & \texttt{00445addbde8} & new\\[4pt]
N2 & \nolinkurl{pa194_ax11h_research_2026-11-14.zip} & \texttt{93ee082ef035} & revised\\[4pt]
N3 & \nolinkurl{pa194_ax11qh4_research_2026-11-14.zip} & \texttt{a21374a58dad} & new\\[4pt]
N4 & \nolinkurl{pa194_frontier_2026-11-14.zip} & \texttt{da53ca17c761} & new\\[4pt]
N5 & \nolinkurl{pa194gh_research_2026-11-14.zip} & \texttt{876db9229518} & new\\[4pt]
\bottomrule
\end{tabular}
\end{center}
Full SHA-256 hashes, exact Git blob identities and prior-version identities are in \nolinkurl{late_snapshot/archive_manifest.json}.


AX11G combines the explicit generalization binder with the four-field,
eleven-entry AX protocol. GH uses the explicit binder with cached-head
Shift; GHS gives its symmetric physical comparison. AX11QH4L and
AX11QH4R reassociate the four-increment constructor in the AX11QH
source program. The two larger incoming surveys repeat AX11G or GH;
the table identities prevent those repeated deliveries from increasing
the graph count.

\begingroup
\footnotesize
\setlength{\tabcolsep}{2pt}
\setlength{\LTleft}{0pt}
\setlength{\LTright}{0pt}
\renewcommand{\arraystretch}{1.1}
\begin{longtable}{@{}>{\raggedright\arraybackslash}p{.07\linewidth}>{\raggedright\arraybackslash}p{.40\linewidth}>{\raggedright\arraybackslash}p{.24\linewidth}>{\raggedright\arraybackslash}p{.22\linewidth}@{}}
\caption{The 16 late-added complete 194-state graph classes}\label{tab:late-classes}\\
\toprule
ID & Graph or search input & Reported status & Class / packet\\
\midrule
\endfirsthead
\multicolumn{4}{@{}l}{\textit{Table \thetable\ continued}}\\
\toprule
ID & Graph or search input & Reported status & Class / packet\\
\midrule
\endhead
\midrule
\multicolumn{4}{r@{}}{\textit{Continued on the next page}}\\
\endfoot
\bottomrule
\endlastfoot
\textbf{L01} & PA194AX11G & principal construction & \texttt{009dba0abc0b}\newline N1, N4\\[4pt]
\textbf{L02} & PA194AX11QH4L & principal construction & \texttt{004bb0473811}\newline N3\\[4pt]
\textbf{L03} & PA194AX11QH4R & principal construction & \texttt{00c276e53e9d}\newline N3\\[4pt]
\textbf{L04} & PA194GH & principal construction & \texttt{008145bc2238}\newline N4, N5\\[4pt]
\textbf{L05} & PA194GHS & principal construction & \texttt{015b0d694cc3}\newline N5\\[4pt]
\textbf{L06} & B/4L search input & research input & \texttt{0059cf22b1db}\newline N3\\[4pt]
\textbf{L07} & AX11QH4L/B search input & research input & \texttt{012e3c3c8501}\newline N3\\[4pt]
\textbf{L08} & AX11QH4R/B search input & research input & \texttt{006f7123f784}\newline N3\\[4pt]
\textbf{L09} & GH neutral edit 84 & research input & \texttt{007a40590f2c}\newline N5\\[4pt]
\textbf{L10} & GH neutral edit 192 & research input & \texttt{023e4487c49e}\newline N5\\[4pt]
\textbf{L11} & GH neutral edit 193 & research input & \texttt{00d5f848002c}\newline N5\\[4pt]
\textbf{L12} & GH neutral edit 296 & research input & \texttt{02fd2818826b}\newline N5\\[4pt]
\textbf{L13} & GH neutral edit 297 & research input & \texttt{0007f5efab7b}\newline N5\\[4pt]
\textbf{L14} & GH neutral edit 298 & research input & \texttt{0262f6022a33}\newline N5\\[4pt]
\textbf{L15} & GH neutral edit 299 & research input & \texttt{00a40ded5b84}\newline N5\\[4pt]
\textbf{L16} & GH neutral edit 300 & research input & \texttt{023dd7565957}\newline N5\\[4pt]
\end{longtable}
\endgroup
The class hashes are 12-digit prefixes of full table-class identifiers, not literal-file hashes. Exact files and all source occurrences are in \nolinkurl{late_snapshot/delta_classes.json}.


\subsection{Novelty corrections and the five released loop constructions}

The table initially constructed under the name PA194GAX11 is
\emph{byte-identical} to the newly delivered PA194AX11G. Both dense files
have SHA-256
\begin{quote}\small
\nolinkurl{7e1c6442816e5df05e961f2057719252a1976c9a195a786978bf1e2d3c8e466d}.
\end{quote}
Its novelty claim is therefore withdrawn against the later snapshot;
it is retained as an inherited comparison artifact.

The original four loop constructions---PA194LC, PA194LC\_A,
PA194LC\_B and PA194LC\_C---remain new against all 110 graphs. The
additional PA194LC\_AX11H table is also new; its dense SHA-256 is
\begin{quote}\small
\nolinkurl{5f10f071e7118ff33aba822e05bae0331b4c78508656c3e945c31b0c6788f9c5}.
\end{quote}
All five are mutually distinct under the same stronger structural test:
arbitrary working-state names, any spanning initial root, left/right
reflection, global symbol exchange, and ignored terminal write/move
decorations. Complete canonical tuples are compared; hashes merely
identify the inputs and receipts. Structural novelty remains separate
from source preservation and arithmetic calibration.

\subsection{Changes in proof status and the evidence boundary}

The
\href{https://github.com/VladimirReshetnikov/ProveIt/blob/603c58ca3bd6bf3222fcfc4e308e6b079f65d573/Computability/BusyBeaver/Research/PA/README.md}{later authoritative PA README}
retains the open status of the primary machine's two absolute
machine-specific directions. Its new formal progress includes a
decomposition of the 58\,154-step startup into the complete source frame,
the first seeded zero-branch instruction reaching site 27 after another
462 transitions, and a finite source-formula translator agreeing with
an independent structural reference on nonstandard valid formula codes.
Complete model-length physical/source correspondence and the internal
proof compiler remain open there.

Separately, the revised AX11H packet now supplies a written mathematical
claim of the absolute iPRA nonhalting/Con(PA) equivalence for its two
literal tables. That claim is stronger than the earlier AX11H delivery;
it is recorded as new attributed proof evidence, not inferred from its
finite execution receipts or from the graph census. Its mathematical
use is discussed separately in this report. AX11G asserts an ordinary
relative iPRA preservation proof, while the GH/GHS and frontier packets
retain an explicit conditional absolute calibration. These distinctions
prevent the late proof claims from being silently substituted into the
initial pinned comparison.

The folder \nolinkurl{late_snapshot/} supplies the full later census,
the 16 additional frozen representatives, full ZIP identities, exact
top-level READMEs, the official status diff and the completed novelty
receipt. From the release root,
\begin{quote}\small\ttfamily
python3 -B late\_snapshot/check\_novelty.py
\end{quote}
repeats the 110-reference comparison offline and checks that GAX11 is
the documented duplicate. The five candidate tables are frozen inputs
to this check; their mathematical correctness is addressed by their
construction and proof sections. The complete later Git-object census
can also be regenerated with \nolinkurl{late_snapshot/rebuild_census.py}.


\section{Searches for a smaller machine}
\label{sec:search}

\subsection{Specified loop families and cross-compositions}

The loop search enumerates every binary association of $k$ ordered
increments followed by a continuation of the following two forms:
$(A,c_2)$, where $A$ contains all $k$ increments, and
$(A,(B,c_3))$, where $A,B$ partition the ordered increments into two
nonempty groups. It tests $4\le k\le9$. Existing literal readers are
reused when their action pairs coincide; missing internal readers are
allocated. The score is the raw control count minus a maximum matching
of the specified disjoint raw-action folding opportunities.

\begin{center}
\begin{tabular}{@{}rrr@{}}
\toprule
Coefficient $k$ & Specified trees & Minimum projected controls\\
\midrule
4 & 10 & 194\\
5 & 28 & 194\\
6 & 84 & 194\\
7 & 264 & 195\\
8 & 858 & 195\\
9 & 2,860 & 195\\
\midrule
Total & 4,104 & 194\\
\bottomrule
\end{tabular}
\end{center}

For coefficient four all ten source graphs are explicitly checked against
the primary. Seven receive projected score 194; three are already present
in the input corpus, while the four promoted new graphs are independently
verified above. Changing $k$ changes the numerical coding, so the higher
coefficient scores are only discovery data until a recoding theorem and
complete physical acceptance checks are supplied. No such candidate
improved the state count.

The three-row LC splice was also applied to 96 distinct frozen compatible
native tables with 194--205 controls, including raw parents, folded
tables and rejected diagnostics. Recomputing the same matching score
gave no value below 194: 73 scored 194, twelve 195, six 196, four 197 and
one 201. The local fold replay is recorded separately. A score alone
does not promote a changed logical design to a consistency machine.
Both enumerations are reproducible from the archive without a repository
checkout or a network connection.

The GAX11 grammar search additionally records 23,500,000 completed
iteration slots across its original core and four loop-core compositions.
All completed runs reached best effective-reader score 60 and produced
no verified sub-194 machine. Three runs with corrected source-address
metadata were interrupted; their logs record at least 4,300,000 further
slots, also without an improved best score. The ledger distinguishes
requested budgets, actual completed loops, skipped proposals and
interrupted checkpoints. The scores are heuristic discovery objectives,
with conservative helper metadata in the first combined runs; they are
not exact lower bounds on physical implementations.

\subsection{Incoming-edge elimination}

The low-level deletion search is more general than merging two literal
control rows. The removed control must differ from the catalogue-selected
start. For each eligible proposed removed control, every incoming row may
be independently redirected, with its write bit and move direction also
changed. A candidate detour must match an old positive-length prefix in
control, head displacement and the entire finite tape window, for every
window in a reachable overapproximation. The detour may use existing
unchanged rows, but may not execute the removed control or another row
being edited. This restriction makes the independently selected
incoming-row replacements compatible.

All 94 previous complete 194-state graphs were screened. The 89 filed
graphs used a 13-cell window, old prefixes of at most 12 transitions,
and replacement paths of at most 128 transitions. The five additional
incoming graphs used a 17-cell window and 256-transition replacement
bound. All runs completed without a deletion. Deeper checks of the
primary and two of the new loop graphs used 19-cell windows, old prefixes
up to 48 and replacements up to 512, again without a deletion.

The finite window abstraction allows either value for a newly exposed
boundary cell. It therefore contains every concrete blank execution from
that catalogue-selected start, but
can admit unreachable windows. This conservatism can reject a valid
optimization that depends on a stronger source invariant. Longer
detours, different encodings, simultaneous multi-control removal,
additional helper controls and startup redesign are outside the tested
rule. The negative result is not a minimality theorem.

\subsection{Shorter-scanner experiments}

Combining the new loop associations with inherited root/MP source
rotations, a standalone history reader and a shortened return scanner
produced several nominal 194-state tables. Their resolved source graphs
and local contextual folds passed, but each failed 47 of the 4,780
universal physical source outcomes. Those tables are rejected. A further
search over a repaired reverse helper and AX grammar changes was
interrupted after its last completed checkpoint of 2,260,000 proposals
and approximately 209 seconds. Its best penalized score was 65 against
a target of 59; these are discovery scores, not complete machine counts.
The checkpoint and limits are retained in the research receipts. The
planned 480-second budget was not completed, and no candidate from that
search is promoted in this report.

\subsection{A particularly instructive rejected 193-state table}

A source splice appeared to reuse an existing increment-and-return
subtree and yielded 193 literal working states. It passed 32 local fold
cases and reached a correctly formed startup frame after 58,320 steps.
It nevertheless had 3,734 source instructions, instead of 3,673. Exact
resolution exposed the following positive branch at source instruction
123:
\begin{lstlisting}
123: dec 7   zero -> 130   positive -> 124
124: inc 1 -> 125
125: inc 1 -> 126
126: inc 1 -> 127
127: inc 1 -> 128
128: inc 6 -> 129
129: jump 124
\end{lstlisting}
The repeated six-instruction cycle adds four to register 1 and one to
register 6. It never returns to the register-7 decrement. This is a
definitive source mismatch, not an unconfirmed test timeout. Correctly
implementing this changed source on tape would still implement the wrong
algorithm. The archive labels the table \texttt{REJECTED} and includes
a standalone reproducer of the counterexample.

\section{What transfers from the suggested adjacent projects}

The ZFC work provides useful engineering methods: command-local field
readers, command-order changes, history-sensitive selectors, exact
whole-tape detours, and independent composition checks. The accepted
K215 construction, for example, prices the complete composed table,
checks the source graph, proves the local marker folds, and separately
checks all arbitrary-register runtime cases and exact startup. It also
distinguishes an external correctness theorem from its internal PA
endpoint.\cite{zfc} Those methods informed the decoder, deletion and
scanner investigations here. The numeric savings of two components do
not automatically add: they can consume the same reader or history
marker, and a source rotation can change the context in which a fold
was valid.

The DiophantineConsistency project is relevant to proof structure. It
keeps the exact literal object, the witness domain, the arithmetic
calibration, and native versus translated consistency sentences separate.
Its reported small equations measure variables, polynomial operations
and degree, often with substantial fixed coefficients.\cite{diophantine}
A small polynomial description does not directly give a small blank-input
binary TM. A construction using it would still need tuple enumeration,
arithmetic evaluation, fixed-constant generation, an exact state census,
and a PA-specific calibration. The cited maintained problem is Con(ZFC),
so using it without a theory-specific redesign would not answer the
present Con(PA) question.

ConciseIPRA offers small proof-checking calculi with exact ordinary
bi-interpretations. Its scalar language uses zero, successor and an
iterator, and builds logical operations from a compact equational
core.\cite{concise-ipra} Its checker-size metrics are normalized Python
tokens, AST nodes and bytes. Replacing the PA machine's numerical
calculus by such a checker would require a separate correctness and
equiconsistency argument for the intended PA proof predicate, together
with literal costs for binders, substitution, freshness, iteration and
the blank loader. A small presentation of iPRA is not, by itself, a
calibrated recognizer of PA refutations. This route remains a worthwhile
architectural redesign, but it supplies no justified state subtraction
from the present table.

The practical implication is specific: future reductions should account
jointly for logical head order, command field order, repeated source
subtrees, tape history markers and initialization. Counting an isolated
smaller formula or source tree is insufficient. The source-preserving LC
change is useful because its entire logical calibration obligation
reduces to that of the unchanged primary program, leaving only physical
implementation and arithmetic synchronization to prove.

% New fragment: PA194GAX11. The historical PA194G text is reference_g_proof.tex.
% Requires amsmath, amssymb and url (or xurl). This fragment asserts a relative result only.
\section{The GAX11 composition: a changed command interface}

The independently reproduced PA194GAX11 combines the AX-parameter generalizer
of PA194G with the four-field fallthrough protocol of PA194AX11.
These two ingredients are inherited. The final snapshot contains an
identical table named PA194AX11G, so this composition is retained as a
checked concurrent reproduction and is excluded from the five-new-graph
claim. Its exact source binding and compiler analysis remain useful.
The dense table has 194 working states, start state zero, and terminal
label 194, which is excluded from the count. It starts on a blank binary
two-way tape and uses only left and right moves. Its SHA-256 is
\begin{quote}\small\texttt{7e1c6442816e5df05e961f2057719252a1976c9a195a786978bf1e2d3c8e466d}.\end{quote}

The result below is a source/command relative halting argument. This
reproduction does not complete its primitive-recursive witness
internalization or import an absolute calibration. The revised AX11H
written theorem is used separately for LC\_AX11H in the main report.
The ordinary PA calibration in the historical references remains
identified as an inherited mathematical claim.

\subsection{Exact command semantics}

Let $p(a,b)=(a+b)(a+b+1)/2+b$, write $[a;w]=p(a,w)$, and use
$I(A,B)=4p(A,B)$ and $U(a,B)=4p(B,a)+2$.
The store $(t,w,d,h)$ consists of the current theorem code, Cantor stack,
numeric offset, and cached theorem code. The latter is the \emph{raw}
incoming cache; it is not the offset-adjusted parameter $H$.

Write $h_1,\ldots,h_5,\rho,\Gamma,g$ for the same numerical families as
the primary separated-diagonal calculus. Use the exchanged sixth head
\[
 h_6^Q(D,c,H)=U(D,I(c,H))\to\bigl(U(D,c)\to U(D,H)\bigr).
\]
This is interderivable with the primary sixth head by the established
implicational exchange template, without altering the first head used in
the derivation of explosion. Set $D=d+x$, $H=h+y$, and
\[
\begin{array}{c|l}
j&q_j\\\hline
0&g(H,D,c)\\
1&\Gamma(c)\\
2&h_6^Q(D,c,H)\\
3&h_5(D,H)\\
4&h_4(D,c,H)\\
5&h_3(D,c,H)\\
6&h_2(D,c,H)\\
7&h_1(D,H)\\
8&\rho(D,c,H,d)\\
9&U(c,h).
\end{array}
\]
As in the primary calculus, $g$ is zero if either original numerical
guard comparison fails. The new command effect is exactly
\[
 \mathsf A_{c,x,y,z}(t,w,d,h)
       =(0,[q_0,q_1,\ldots,q_9,t;w],z,0).                 \tag{GAX}
\]
Thus AX saves eleven entries and consumes four parameters. The final
field $z$ is real: fallthrough enters the complete MP parameter reader.
Its zero-input MP comparison returns zero whether $z=0$ or $z>0$,
while setting the offset to $z$. The compiler chooses $z=0$, but the
soundness argument covers every natural $z$.

With $v=d+x$, the remaining commands are
\[
 \mathsf{MP}_x(t,w,d,h)=
 \begin{cases}(v,w,v,0),&I(t,v)=h,\\(0,w,v,0),&I(t,v)\ne h,
 \end{cases}
\]
and
\[
 \mathsf S(t,w,d,h)=
 \begin{cases}\mathsf{halt},&t+h=1,\\
 (h,\operatorname{tl}(w),d,\operatorname{hd}(w)),&t+h\ne1.
 \end{cases}
\]
Truth code zero has a fixed calculus proof: the restricted head satisfies
$\rho(0,0,0,0)=0$. This proof supplies virtual entries when the
canonical Cantor stack is empty.

All these equations use total Cantor projections, including
$[0;0]=\operatorname{hd}(0)=\operatorname{tl}(0)=0$.

\subsection{The command compiler, including its boundary cases}

Let $A_0=\mathsf A_{0,0,0,0}$. Since $q_0,\ldots,q_8$ have residue
zero modulo four and $q_9=U(c,h)\ge2$, direct iteration gives
\begin{align*}
 \mathsf A_{c,x,y,0};\mathsf S^{10}(t,w,d,h)
      &=(q_8,[t;w],0,q_9),\\
 \mathsf A_{c,x,y,0};\mathsf S^{11}(t,w,d,h)
      &=(q_9,w,0,t).
\end{align*}
The displayed iterations are nonhalting for every natural store.
Define
\[
 \mathsf{Push}=A_0;\mathsf S^{10};\mathsf{MP}_0,\qquad
 \mathsf{Pop}=A_0;\mathsf S^{12},\qquad
 \mathsf{Discard}=\mathsf{MP}_0.
\]
They satisfy
\begin{align*}
 \mathsf{Push}(t,w,d,h)&=(0,[t;w],0,0),\\
 \mathsf{Pop}(t,[a;w],d,h)&=(t,w,0,a),\\
 \mathsf{Discard}(t,w,0,h)&=(0,w,0,0).
\end{align*}
Pop preserves even the visible contradiction code $t=1$: its final
test includes a universal code at least two. No free-list assumption is
used, so zero pushes and the empty stack are covered.

From $(t,w,0,0)$, the selector of $q_j$, $0\le j\le8$, is
\[
 \mathsf{Discard};\mathsf A_{c,D,H,0};
           \mathsf S^{j+2};\mathsf{Pop}^{9-j}.
\]
It returns $(q_j,w,0,0)$. For the restricted head with an actual
parameter $u\le D$, use
\[
 \mathsf{Discard};\mathsf{MP}_u;
 \mathsf A_{c,D-u,H,0};\mathsf S^{10};\mathsf{Pop}.
\]
The important new generalization word is
\[
 \boxed{\mathsf{Gen}_i=
       \mathsf{Push};\mathsf S;\mathsf A_{i,0,0,0};\mathsf S^{11}.}
\]
It maps $(a,w,0,0)$ to $(U(i,a),w,0,0)$ for all $a,w,i$.
No offset-setting command is needed because the binder is the explicit
first AX parameter. The intermediate Shift sees zero in both visible
positions, and the stored arbitrary theorem $a$ is used as the raw cache
of AX. This includes $a=1$ and $i=0$.

On a store $(0,[a,I(a,b),0;w],0,0)$ the suffix
\[
       \mathsf S;\mathsf S;\mathsf{MP}_b;\mathsf{Pop}
\]
halts if $a=1$ and otherwise returns $(b,w,0,0)$.
The reserved zero beneath the premises is necessary because successful
MP does not push a dummy zero. The case $b=1$ returns the contradiction
without losing it; a final plain Shift then halts. If $a=b=w=0$, the
same calculation holds using total projections despite all zero aliases.

Consequently, recursively compile a numerical proof $e$ into a word
$W_e$. Use the selectors for its axiom leaves, append $\mathsf{Gen}_i$
at a generalization node, and for an MP node with major proof $e_1$ of
$I(a,b)$ and minor proof $e_0$ of $a$ use
\[
 \mathsf{Discard};\mathsf{Push};W_{e_1};\mathsf{Push};
 W_{e_0};\mathsf{Push};\mathsf S;\mathsf S;
                       \mathsf{MP}_b;\mathsf{Pop}.
\]
Induction on the proof tree proves the precise invariant: from every
$(t,w,0,0)$ the word either halts or returns
$(\operatorname{concl}(e),w,0,0)$. In particular, when $e$ proves one,
$A_0;W_e;\mathsf S$ halts from \emph{every} carried store.

Conversely, annotate the current top, cached theorem, and canonical
live stack entries with numerical-calculus derivations. Every head of
(GAX) is a theorem; the new head $U(c,h)$ is obtained by generalizing
the cached derivation. The extra MP field returns truth and has no
theorem obligation on its numeric offset. MP and Shift preserve the
annotations. If $t+h=1$, one annotated theorem is one, yielding a
refutation. This argument never treats $t=h=1$ as a halting condition.

\subsection{Scheduling and the relative arithmetic statement}

Let $e(\epsilon)=0$, $e(C::w)=1+b_C(e(w))$, where
\begin{align*}
 b_{\mathsf S}(r)&=p(0,r),\\
 b_{\mathsf{MP}_x}(r)&=p(1,p(x,r)),\\
 b_{\mathsf A_{c,x,y,z}}(r)&=p(c+2,p(y,p(x,p(z,r)))).
\end{align*}
The initial raw scheduler candidate for a nonempty word $C::w$ is
$b_C(e(w))$; the value $1+b_C(e(w))$ is its encoded nonempty
continuation, which is decremented before command decoding. This
distinguishes the root candidate from nested continuations.

The extra field must be present; otherwise AX consumes part of the
following command as $z$. Total inverse identities decode these words
even when fields and tails are zero. Each positive continuation decreases
strictly after a command. Prefixing any all-store halting suffix with
$\mathsf A_{k,0,0,0}$ gives an unbounded family of candidate codes, so
the unchanged increasing scheduler reaches a halting candidate unless
it halts earlier. This uses the all-store property, not a reset between
candidate words.

At the source level these two transformations show halting equivalence
with the numerical calculus, and hence with the PA194G source. They use
finite proof recursion, arithmetic command calculations, and witnessed
finite scheduling. Together with the source/tape simulations they yield
a relative transfer of the blank nonhalting sentences. Expressing those
finite transformations as primitive-recursive witness functions and
checking their correctness is the weak-arithmetic formalization task;
no end-to-end Lean theorem for this composition is supplied here.

In particular, the logical implication from an independently established
baseline iPRA calibration and an iPRA proof of this relative transfer to
the corresponding iPRA calibration of GAX11 is immediate. The artifact
does not silently replace either premise by a standard-Nat execution
test. The stronger absolute statement
$\mathrm{iPRA}\vdash\neg\mathrm{Halt}(\mathrm{GAX11})
\leftrightarrow\mathrm{Con}(\mathrm{PA})$
is therefore not recorded as a completed theorem of this release.

\subsection{Literal source binding and executed evidence}

Relative to the PA194G source, the sole new source splice replaces its
AX return at ordinal 3553 by a nineteen-instruction push. Its normal exit
enters the complete MP reader. All 3653 other instructions, including
their resolved successors, agree under the relocation
$i\mapsto i$ for $i\le3553$ and $i\mapsto i+18$ otherwise.
No external source edge enters the inserted block's interior.
Thus the new source has 3672 primitive sites.

The table is compiled from 134 fixed controls and 67 AX readers, followed
by seven simultaneous contextual folds. All 28 local fold cases pass in
the final simultaneous table. The construction checker verifies all 4779
source outcomes with arbitrary natural registers; its recorded workload
is 1483363 symbolic operations with maximum 470 per outcome. A complete
58154-step blank initialization reaches source site 25 with the usual
logical register vector $(2,0,\ldots,0)$ and 26 unit spare blocks.
These are implementation checks, not finite-run evidence for consistency.

The independent literal-tape replay also completed successfully. It
checked all 4,779 outcomes using 1,083,079 symbolic operations, with maximum
379 per outcome and minimum positive duration constant 3. It derived
four generic list lemmas from literal rows, compared two independently
resolved source graphs, and independently checked the complete
58,154-transition startup. Its table-bound receipt is retained under
\path{additional/PA194GAX11/evidence_full/independent/PA194GAX11/}.


\section{Reproduction, provenance and delivery contents}

\subsection{Portable replay}

The main build and symbolic verification use Python 3.10 or later and
its standard library, with assertions enabled. Run from the archive root:
\begin{lstlisting}
python3 reproduce.py
python3 reproduce.py --physical
python3 research/enumerate_loop_trees.py
python3 research/patch_loop_corpus.py
python3 research/reject_false193.py
\end{lstlisting}
The first command reconstructs all five new machines in isolated
temporary directories and compares the generated artifacts to the release.
It binds the recorded independent certificates to the exact tables and
contracts, reruns source, startup and affine-clock checks, and verifies
novelty against all 110 prior graphs. The second additionally regenerates
the physical contracts for both parents and all five new tables, using
external copies. It can take several minutes. Default replay explicitly
records that it did not regenerate universal contracts. Both modes write
new evidence outside the frozen release; an optional \texttt{--out PATH}
selects that directory. Neither verifies the full imported arithmetic
derivation or runs a proof-assistant kernel.

The principal machine also has its own combined verifier in
\path{additional/PA194LC_AX11H/verify_all.py}. The retained concurrent
GAX11 reproduction has an isolated verifier in
\path{additional/PA194GAX11/verify_all.py}. The bounded incoming-edge
search in \path{research/incoming_edge_search/} additionally requires a
C++17 compiler. Its catalogue pins all 94 frozen inputs. The full Git
corpus census can be regenerated with the archived census tool and a
checkout containing the pinned Git blobs; reading the supplied census
and replaying novelty against its frozen representatives needs no Git
checkout. The report and the two complete census appendices compile
with standard pdfLaTeX packages.

\subsection{Artifact map}

\begin{center}\small
\begin{tabularx}{\textwidth}{@{}p{.35\textwidth}X@{}}
\toprule
Path & Contents\\
\midrule
\path{report/} & This manuscript, all included \texttt{.tex} fragments, compiled PDF, full LC\_AX11H and LC tables, and census appendices.\\
\path{machines/} & Four final dense/native tables, source trees, complete resolved instructions, folds, startup records and independent contracts.\\
\path{reference/} & Frozen original-primary data and the unchanged revised AX11H calibration ZIP.\\
\path{construction/} & Deterministic standalone builder and exact four-variant specifications.\\
\path{verification/} & Independent literal-row checker, source resolvers, finite replay tools, polynomial-envelope receipts and proof review.\\
\path{comparison/} & 82-packet ledger, 94-graph annotations, full census, frozen representative tables and canonical-comparison tools.\\
\path{research/} & Bounded searches, frozen inputs, negative results and the explicitly rejected 193-state diagnostic.\\
\path{additional/PA194LC_AX11H/} & Principal new machine, frozen AX11H parent, standalone relative proof, source identity, complete physical evidence and verifier.\\
\path{late_snapshot/} & Final census, 16 added prior graphs, five incoming ZIP identities and the 110-reference novelty replay.\\
\path{additional/PA194GAX11/} & Concurrent AX11G reproduction, source/command argument, isolated build and physical checks.\\
\path{licenses/} & Original repository and upstream source notices.\\
\bottomrule
\end{tabularx}
\end{center}

\subsection{Provenance and what was not re-proved}

The primary machine, its numerical calculus, the register engine,
source language, existing contextual folds and prior logical
transformations are inherited research. This report adds five loop
continuation compositions, their deterministic artifact bindings, the
independent replays reported above, the relative clock proofs, the
explicitly attributed AX11H absolute corollary, the two-snapshot corpus
audit and the scoped searches. GAX11 is retained as a concurrent
reproduction after the final comparison identified an identical table.
The original MIT-style notices are retained. A successful current check
does not turn every inherited logical assertion into a newly verified
theorem; the census annotations identify historical evidence separately.

The archive supplies explicit relative iPRA proofs for all five new
source-preserving machines. For PA194LC\_AX11H it also supplies the
absolute corollary importing the revised AX11H written theorem, with
that exact upstream archive and the review attached. The original-primary
four-machine family retains its separately stated relative status.
No new end-to-end Lean or Rocq kernel build is claimed.

\begin{thebibliography}{9}
\bibitem{pa-readme}
Vladimir Reshetnikov and ProveIt contributors.
\emph{PA consistency machines}, authoritative status and machine bindings.
Pinned revision \texttt{32d0d8b7d93a}.
\href{\repourl/Computability/BusyBeaver/Research/PA/README.md}{Repository README}.

\bibitem{machines}
ProveIt contributors. \emph{Delivered PA research packets and machine census}.
Same pinned revision.
\href{\repourl/Computability/BusyBeaver/Research/PA/docs/README.md}{Packet index};
\href{\repourl/Computability/BusyBeaver/Research/PA/docs/MACHINES.md}{Machine census};
\href{\repourl/docs/incoming}{Incoming archives}.
The companion CSVs preserve exact member locators and hashes.

\bibitem{zfc}
ProveIt contributors. \emph{Smaller machines for Con(ZFC)} and
\emph{Ordinary proof of the combined215 machine}. Same pinned revision.
\href{\repourl/Computability/BusyBeaver/Research/ZFC/README.md}{Project status};
\href{\repourl/Computability/BusyBeaver/Research/ZFC/candidate-k215/proof.md}{K215 composition proof}.

\bibitem{diophantine}
ProveIt contributors. \emph{A Diophantine equation for syntactic consistency
of ZFC}. Same pinned revision.
\href{\repourl/Computability/DiophantineConsistency/README.md}{Project README}.

\bibitem{concise-ipra}
ProveIt contributors. \emph{Concise presentations of intuitionistic
primitive recursive arithmetic}. Same pinned revision.
\href{\repourl/Logic/PrimitiveRecursiveArithmetic/ConciseIPRA/README.md}{Project README and article index}.

\bibitem{latest-ax11h}
ProveIt incoming research. \emph{PA194AX11H and PA194QAX11H}, revised
complete written calibration packet, at commit
\texttt{603c58ca3bd6bf3222fcfc4e308e6b079f65d573}.
\href{https://github.com/VladimirReshetnikov/ProveIt/blob/603c58ca3bd6bf3222fcfc4e308e6b079f65d573/docs/incoming/pa194_ax11h_research_2026-11-14.zip}{Exact upstream archive}.
See especially \path{proof/consistency_equivalence.tex},
\path{proof/forward_calibration_completion.tex},
\path{proof/foundation_paeq_completion.tex},
\path{proof/term_leaf_completion.tex},
\path{proof/physical_ipra_audit.tex}, and
\path{proof/source_command_completion.tex}.
An unchanged copy is included in this release.

\bibitem{ipra-transfer}
ProveIt contributors. \emph{Transfer from I$\Sigma_1$}.
Same pinned revision.
\href{\repourl/Computability/BusyBeaver/Research/PA/isigma1-transfer/README.md}{Generic transfer development}.
This supplies a conditional transfer method and no table-specific calibration.
\end{thebibliography}

\clearpage
\appendix
\section{Complete PA194LC\_AX11H transition table}
\label{app:lc-ax11h-table}
This is the complete dense table. Start in control 0 on an all-zero
binary tape; control 194 is the sole terminal and is excluded from the
194 working states. An action $wDq$ writes $w$, moves in direction $D$
and enters $q$. Each group below repeats the columns control/read-0/read-1.
The machine-readable file is \path{additional/PA194LC_AX11H/machine.tm}.

\begin{center}
\scriptsize\ttfamily
\setlength{\tabcolsep}{5pt}
\renewcommand{\arraystretch}{1.15}
\begin{tabular}{@{}rrr@{\hspace{11pt}}rrr@{\hspace{11pt}}rrr@{\hspace{11pt}}rrr@{}}
\toprule
q & 0 & 1 & q & 0 & 1 & q & 0 & 1 & q & 0 & 1 \\
\midrule
0 & 1R1 & 1R2 & 49 & 0R68 & 0R69 & 98 & 0R111 & 1R75 & 147 & 0R149 & 1R158 \\
1 & 1R3 & 1R1 & 50 & 0R70 & 0L48 & 99 & 0L63 & 0R109 & 148 & 0R159 & 1R88 \\
2 & 0R0 & 1R2 & 51 & 0R64 & 1R53 & 100 & 0L62 & 1R113 & 149 & 0R97 & 1R160 \\
3 & 0R4 & 0R1 & 52 & 0L71 & 0L15 & 101 & 0L90 & 1R114 & 150 & 0R161 & 1R162 \\
4 & 1R5 & 1L6 & 53 & 0R72 & 0L10 & 102 & 0R115 & 1R116 & 151 & 0R163 & 1R97 \\
5 & 1L7 & 1L8 & 54 & 0R73 & 1R74 & 103 & 0R69 & 1R103 & 152 & 0R164 & 1R165 \\
6 & 0L7 & 1L6 & 55 & 0R75 & 1R76 & 104 & 0R117 & 1R95 & 153 & 0R166 & 1R157 \\
7 & 0L9 & 1L6 & 56 & 0R77 & 1R78 & 105 & 0R117 & 1R67 & 154 & 0R167 & 1R156 \\
8 & 0L10 & 0L11 & 57 & 0R79 & 1R80 & 106 & 0R118 & 0L14 & 155 & 0R135 & 1R130 \\
9 & 1R12 & 0L9 & 58 & 0L81 & 0R1 & 107 & 0R119 & 1L71 & 156 & 0R147 & 1R41 \\
10 & 0L9 & 0L10 & 59 & 0R82 & 1R59 & 108 & 0R120 & 1R121 & 157 & 0R134 & 1R168 \\
11 & 1L8 & 1L11 & 60 & 0L73 & 1L46 & 109 & 0R14 & 1R109 & 158 & 0R41 & 1L169 \\
12 & 0R12 & 0R13 & 61 & 0R83 & 1R84 & 110 & 0R66 & 1R110 & 159 & 0L170 & 1R171 \\
13 & 0L14 & 0R13 & 62 & 0R54 & 1R85 & 111 & 0R122 & 1L35 & 160 & 0R107 & 1R85 \\
14 & 1L15 & 0L15 & 63 & 0L86 & 1R87 & 112 & 0R123 & 1R124 & 161 & 0R172 & 1R170 \\
15 & 0L15 & 1L16 & 64 & 0R58 & 0R82 & 113 & 0R88 & 0L14 & 162 & 0R173 & 1R174 \\
16 & 0L17 & 1L16 & 65 & 0R88 & 1R30 & 114 & 0R90 & 1L95 & 163 & 0R175 & 1R176 \\
17 & 0L18 & 1L16 & 66 & 0R89 & 1R66 & 115 & 0R88 & 1R88 & 164 & 0R177 & 1R178 \\
18 & 0L19 & 1L16 & 67 & 0L90 & 1L33 & 116 & 0R115 & 1L37 & 165 & 0L179 & 1L180 \\
19 & 0L20 & 1L16 & 68 & 0L81 & 0R2 & 117 & 0L63 & 0R103 & 166 & 0R149 & 1R140 \\
20 & 0L21 & 1L16 & 69 & 0R91 & 1R69 & 118 & 0R68 & 0R91 & 167 & 0R165 & 1R181 \\
21 & 0L22 & 1L16 & 70 & 0R72 & 1R92 & 119 & 0R58 & 0R91 & 168 & 0L130 & 1R88 \\
22 & 0L23 & 1L16 & 71 & 0R46 & 1R93 & 120 & 0L86 & 1R125 & 169 & 0R182 & 1R135 \\
23 & 0R23 & 1R24 & 72 & 0R42 & 0L13 & 121 & 0R126 & 1R127 & 170 & 0R183 & 1R145 \\
24 & 0R25 & 1R26 & 73 & 0R94 & 1R37 & 122 & 0R58 & 0R110 & 171 & 0R158 & 1R130 \\
25 & 0R27 & 1R28 & 74 & 0R30 & 1R95 & 123 & 0R99 & 1R113 & 172 & 0R139 & 1R165 \\
26 & 0R29 & 1R30 & 75 & 0R42 & 1R96 & 124 & 0R79 & 1R78 & 173 & 0R158 & 1R184 \\
27 & 0R31 & 1R194 & 76 & 0R97 & 1R98 & 125 & 0R128 & 1R129 & 174 & 0R185 & 1R107 \\
28 & 0L32 & 1R33 & 77 & 0R99 & 1L100 & 126 & 0R32 & 1R36 & 175 & 0R31 & 1R52 \\
29 & 0R34 & 1R35 & 78 & 0R30 & 1R77 & 127 & 0R31 & 1R98 & 176 & 0R186 & 1R175 \\
30 & 0L36 & 1L37 & 79 & 0R93 & 1R101 & 128 & 0R130 & 1R131 & 177 & 0R134 & 1R142 \\
31 & 0R38 & 1R39 & 80 & 0R30 & 1R102 & 129 & 0R132 & 1R133 & 178 & 0R32 & 1R187 \\
32 & 0R40 & 1R41 & 81 & 1R81 & 1L15 & 130 & 0R134 & 1R135 & 179 & 0R144 & 1R188 \\
33 & 0R42 & 1R43 & 82 & 0R103 & 1R82 & 131 & 0R136 & 1R137 & 180 & 0R189 & 1R190 \\
34 & 0R44 & 1R45 & 83 & 0R94 & 0L7 & 132 & 0R138 & 1R139 & 181 & 0R41 & 1R191 \\
35 & 0R46 & 1R33 & 84 & 0R104 & 1R105 & 133 & 0R140 & 1R141 & 182 & 0R168 & 1R41 \\
36 & 0R47 & 1R48 & 85 & 0R42 & 1R106 & 134 & 0R97 & 1R35 & 183 & 0R192 & 1R88 \\
37 & 0R49 & 1R50 & 86 & 0R107 & 1R48 & 135 & 0R142 & 1R115 & 184 & 0R153 & 1R146 \\
38 & 0R51 & 1R42 & 87 & 0R108 & 0L72 & 136 & 0R143 & 1R144 & 185 & 0R165 & 1R180 \\
39 & 0R52 & 1R53 & 88 & 0R68 & 0R66 & 137 & 0R145 & 1R146 & 186 & 0R92 & 1R111 \\
40 & 0R54 & 1R55 & 89 & 0R109 & 1R89 & 138 & 0R134 & 1R147 & 187 & 0R193 & 1R159 \\
41 & 0R56 & 1R57 & 90 & 0R49 & 1R100 & 139 & 0R88 & 1R148 & 188 & 0R169 & 1R154 \\
42 & 0R58 & 0R59 & 91 & 0R110 & 1R91 & 140 & 0R149 & 1L131 & 189 & 0R158 & 1R192 \\
43 & 0R60 & 0L14 & 92 & 0R30 & 1R48 & 141 & 0R150 & 1R151 & 190 & 0R179 & 1R182 \\
44 & 0R61 & 1R62 & 93 & 0R47 & 1R60 & 142 & 0R76 & 1R41 & 191 & 0R182 & 1R41 \\
45 & 0R63 & 1R52 & 94 & 0L63 & 0R89 & 143 & 0R28 & 1R152 & 192 & 0R149 & 1R168 \\
46 & 0R64 & 1R65 & 95 & 0R94 & 1L111 & 144 & 0L153 & 1L154 & 193 & 0R76 & 1R144 \\
47 & 0R58 & 0R66 & 96 & 0R101 & 0L14 & 145 & 0L129 & 1R155 &  &  &  \\
48 & 0R67 & 0L14 & 97 & 0R36 & 1R112 & 146 & 0R156 & 1R157 &  &  &  \\
\bottomrule
\end{tabular}
\end{center}

\clearpage
\section{Complete PA194LC transition table}
\label{app:lc-table}
This is the complete dense table. Start in control 0 on an all-zero
binary tape; control 194 is the sole terminal and is excluded from the
194 working states. An action $wDq$ writes $w$, moves in direction $D$
and enters $q$. Each group below repeats the columns control/read-0/read-1.
The machine-readable file is \path{machines/PA194LC/machine.tm}.

\begin{center}
\scriptsize\ttfamily
\setlength{\tabcolsep}{5pt}
\renewcommand{\arraystretch}{1.15}
\begin{tabular}{@{}rrr@{\hspace{11pt}}rrr@{\hspace{11pt}}rrr@{\hspace{11pt}}rrr@{}}
\toprule
q & 0 & 1 & q & 0 & 1 & q & 0 & 1 & q & 0 & 1 \\
\midrule
0 & 1R1 & 1R2 & 49 & 0R69 & 0L14 & 98 & 0R87 & 1L115 & 147 & 0R149 & 1R158 \\
1 & 1R3 & 1R1 & 50 & 0R70 & 0R71 & 99 & 0R104 & 0L14 & 148 & 0R159 & 1R121 \\
2 & 0R0 & 1R2 & 51 & 0R72 & 0L49 & 100 & 0R36 & 1R116 & 149 & 0R100 & 1R160 \\
3 & 0R4 & 0R1 & 52 & 0R73 & 1R54 & 101 & 0R115 & 1R78 & 150 & 0R161 & 1R162 \\
4 & 1R5 & 1L6 & 53 & 0L74 & 0L15 & 102 & 0L46 & 0R113 & 151 & 0R163 & 0L10 \\
5 & 1L7 & 1L8 & 54 & 0R75 & 0L10 & 103 & 0L45 & 1R117 & 152 & 0R164 & 1R165 \\
6 & 0L7 & 1L6 & 55 & 0R76 & 1R77 & 104 & 0L94 & 1R118 & 153 & 0R166 & 1R157 \\
7 & 0L9 & 1L6 & 56 & 0R78 & 1R79 & 105 & 0R119 & 1R120 & 154 & 0R167 & 1R156 \\
8 & 0L10 & 0L11 & 57 & 0R80 & 1R81 & 106 & 0R71 & 1R106 & 155 & 0R135 & 1R130 \\
9 & 1R12 & 0L9 & 58 & 0R82 & 1R83 & 107 & 0R121 & 1R30 & 156 & 0R147 & 1R41 \\
10 & 0L9 & 0L10 & 59 & 0L84 & 0R1 & 108 & 0L46 & 0R106 & 157 & 0R134 & 1R168 \\
11 & 1L8 & 1L11 & 60 & 0R85 & 1R60 & 109 & 0R70 & 0R95 & 158 & 0R41 & 1L169 \\
12 & 0R12 & 0R13 & 61 & 0L76 & 1L86 & 110 & 0R59 & 0R95 & 159 & 0L170 & 1R171 \\
13 & 0L14 & 0R13 & 62 & 0R87 & 0L7 & 111 & 0L65 & 1R122 & 160 & 0R91 & 1R64 \\
14 & 1L15 & 0L15 & 63 & 0R88 & 1R89 & 112 & 0R123 & 1R124 & 161 & 0R172 & 1R170 \\
15 & 0L15 & 1L16 & 64 & 0R42 & 1R90 & 113 & 0R14 & 1R113 & 162 & 0R173 & 1R174 \\
16 & 0L17 & 1L16 & 65 & 0R91 & 1R49 & 114 & 0R68 & 1R114 & 163 & 0R175 & 1R176 \\
17 & 0L18 & 1L16 & 66 & 0R92 & 0L75 & 115 & 0R125 & 1L67 & 164 & 0R177 & 1R178 \\
18 & 0L19 & 1L16 & 67 & 0R86 & 1R33 & 116 & 0R126 & 1R127 & 165 & 0L179 & 1L180 \\
19 & 0L20 & 1L16 & 68 & 0R93 & 1R68 & 117 & 0R121 & 0L14 & 166 & 0R149 & 1R140 \\
20 & 0L21 & 1L16 & 69 & 0L94 & 1L33 & 118 & 0R94 & 1L98 & 167 & 0R165 & 1R181 \\
21 & 0L22 & 1L16 & 70 & 0L84 & 0R2 & 119 & 0R121 & 1R121 & 168 & 0L130 & 1R121 \\
22 & 0L23 & 1L16 & 71 & 0R95 & 1R71 & 120 & 0R119 & 1L37 & 169 & 0R182 & 1R135 \\
23 & 0R23 & 1R24 & 72 & 0R75 & 1R96 & 121 & 0R70 & 0R68 & 170 & 0R183 & 1R145 \\
24 & 0R25 & 1R26 & 73 & 0R59 & 0R85 & 122 & 0R128 & 1R129 & 171 & 0R158 & 1R130 \\
25 & 0R27 & 1R28 & 74 & 0R86 & 1R97 & 123 & 0R32 & 1R36 & 172 & 0R139 & 1R165 \\
26 & 0R29 & 1R30 & 75 & 0R42 & 0L13 & 124 & 0R31 & 1R79 & 173 & 0R158 & 1R184 \\
27 & 0R31 & 1R194 & 76 & 0R87 & 1R37 & 125 & 0R59 & 0R114 & 174 & 0R185 & 1R91 \\
28 & 0L32 & 1R33 & 77 & 0R30 & 1R98 & 126 & 0R102 & 1R117 & 175 & 0R31 & 1R53 \\
29 & 0R34 & 1R35 & 78 & 0R42 & 1R99 & 127 & 0R82 & 1R81 & 176 & 0R186 & 1R175 \\
30 & 0L36 & 1L37 & 79 & 0R100 & 1R101 & 128 & 0R130 & 1R131 & 177 & 0R134 & 1R142 \\
31 & 0R38 & 1R39 & 80 & 0R102 & 1L103 & 129 & 0R132 & 1R133 & 178 & 0R32 & 1R187 \\
32 & 0R40 & 1R41 & 81 & 0R30 & 1R80 & 130 & 0R134 & 1R135 & 179 & 0R144 & 1R188 \\
33 & 0R42 & 1R43 & 82 & 0R97 & 1R104 & 131 & 0R136 & 1R137 & 180 & 0R189 & 1R190 \\
34 & 0R44 & 1R45 & 83 & 0R30 & 1R105 & 132 & 0R138 & 1R139 & 181 & 0R41 & 1R191 \\
35 & 0R46 & 1R47 & 84 & 1R84 & 1L15 & 133 & 0R140 & 1R141 & 182 & 0R168 & 1R41 \\
36 & 0R48 & 1R49 & 85 & 0R106 & 1R85 & 134 & 0R100 & 1R67 & 183 & 0R192 & 1R121 \\
37 & 0R50 & 1R51 & 86 & 0R73 & 1R107 & 135 & 0R142 & 1R119 & 184 & 0R153 & 1R146 \\
38 & 0R52 & 1R42 & 87 & 0L46 & 0R93 & 136 & 0R143 & 1R144 & 185 & 0R165 & 1R180 \\
39 & 0R53 & 1R54 & 88 & 0R108 & 1R98 & 137 & 0R145 & 1R146 & 186 & 0R96 & 1R115 \\
40 & 0R55 & 1R56 & 89 & 0R108 & 1R69 & 138 & 0R134 & 1R147 & 187 & 0R193 & 1R159 \\
41 & 0R57 & 1R58 & 90 & 0R109 & 0L14 & 139 & 0R121 & 1R148 & 188 & 0R169 & 1R154 \\
42 & 0R59 & 0R60 & 91 & 0R110 & 1L74 & 140 & 0R149 & 1L131 & 189 & 0R158 & 1R192 \\
43 & 0R61 & 0L14 & 92 & 0R111 & 1R112 & 141 & 0R150 & 1R151 & 190 & 0R179 & 1R182 \\
44 & 0R62 & 1R63 & 93 & 0R113 & 1R93 & 142 & 0R79 & 1R41 & 191 & 0R182 & 1R41 \\
45 & 0R55 & 1R64 & 94 & 0R50 & 1R103 & 143 & 0R28 & 1R152 & 192 & 0R149 & 1R168 \\
46 & 0L65 & 1R66 & 95 & 0R114 & 1R95 & 144 & 0L153 & 1L154 & 193 & 0R79 & 1R144 \\
47 & 0R36 & 1R67 & 96 & 0R30 & 1R49 & 145 & 0L129 & 1R155 &  &  &  \\
48 & 0R59 & 0R68 & 97 & 0R48 & 1R61 & 146 & 0R156 & 1R157 &  &  &  \\
\bottomrule
\end{tabular}
\end{center}

\clearpage
\section{Exact identities of the delivered tables}
The SHA-256 values below refer to the exact UTF-8 table files, including their final newlines. The archive manifest also binds source and certificate files.
\begingroup\small
\begin{longtable}{@{}p{.23\textwidth}p{.72\textwidth}@{}}
\toprule
File & SHA-256\\
\midrule\endhead
LC\_AX11H dense & \nolinkurl{5f10f071e7118ff33aba822e05bae0331b4c78508656c3e945c31b0c6788f9c5}\\[4pt]
LC\_AX11H native & \nolinkurl{e7800298659a5768a8aabfbd0cfd0ea42af241ed2a10b0cc1dbf64ed8aff1112}\\[4pt]

PA194LC dense & \nolinkurl{9cb620ba1b869fb793fd9ba0a99e90dcc85d9a6d62a42f48109743a7fbeda3c8}\\[4pt]
PA194LC native & \nolinkurl{ee09a6e73bb73393d3743795a5c46bcbe1a524b7ce4fd685a511562bd9eb5642}\\[4pt]
PA194LC\_A dense & \nolinkurl{382663bad011e958b51221a589e47880d434172d2f3417d2f3e940acda33cbdf}\\[4pt]
PA194LC\_A native & \nolinkurl{396d4425f11a31184abc7500d326441b073c4035e810d1a1fe13f77e405698fd}\\[4pt]
PA194LC\_B dense & \nolinkurl{4b2b45e97fdd5f8a6e22c4f552dc3f26252c170e85cdb5ddb27ad5d5c580982a}\\[4pt]
PA194LC\_B native & \nolinkurl{70a6cec8eba78f8fd279045c330a849fc5cb5f8dc420a038debf7618e8744de6}\\[4pt]
PA194LC\_C dense & \nolinkurl{7de80d804e5994db6008687421b0c5b4e96c9c87097fb293269406e9009a5753}\\[4pt]
PA194LC\_C native & \nolinkurl{5004d00641875679881c0f8d311386a675f09b18c123acd275640a79c231db3c}\\[4pt]
GAX11 (reproduced) & \nolinkurl{7e1c6442816e5df05e961f2057719252a1976c9a195a786978bf1e2d3c8e466d}\\[4pt]
\bottomrule\end{longtable}\endgroup
The primary baseline identities are given in Section 2. Native identifiers are carried through the exact label arrays; maximum native label is not a state count.

\clearpage
\section{The 94 initial-snapshot 194-state graph classes}

\label{app:frontier194}

This appendix exhausts the complete spanning 194-state graphs recognized in the pinned supplied corpus; it excludes the constructions first delivered in this report. The category is an attributed packet status: P = reported principal construction (51), S = supplementary construction (7), A = research auxiliary (25), D = search seed (4), R = rejected (7). None of these letters certifies a completed absolute iPRA consistency theorem.

The short hash is the first 12 hexadecimal digits of the canonical table-class identifier, not a literal-file hash. All 94 prefixes are distinct. The source column names one packet containing a representative; it is not necessarily the first origin of that design. Packet identifiers are expanded in Appendix~\ref{app:packet-ledger}. The full hashes, literal-file locators and JSON pointers are in \nolinkurl{comparison/frontier194_annotated.csv}; \nolinkurl{report/graph_packet_mapping.csv} binds the short identifiers below to that data.
\begingroup
\footnotesize
\setlength{\tabcolsep}{2pt}
\setlength{\LTleft}{0pt}
\setlength{\LTright}{0pt}
\renewcommand{\arraystretch}{1.12}
\begin{longtable}{@{}>{\raggedright\arraybackslash}p{.16\linewidth}>{\raggedright\arraybackslash}p{.19\linewidth}>{\raggedright\arraybackslash}p{.045\linewidth}>{\raggedright\arraybackslash}p{.065\linewidth}>{\raggedright\arraybackslash}p{.49\linewidth}@{}}
\caption{Complete initial-snapshot 194-state graph census}\label{tab:frontier194}\\
\toprule
ID / class hash & Name & Cat. & Source & Mechanism\\
\midrule
\endfirsthead
\multicolumn{5}{@{}l}{\textit{Table \thetable\ continued}}\\
\toprule
ID / class hash & Name & Cat. & Source & Mechanism\\
\midrule
\endhead
\midrule
\multicolumn{5}{r@{}}{\textit{Continued on the next page}}\\
\endfoot
\bottomrule
\endlastfoot
\hypertarget{graph-G01}{\textbf{G01}}\newline{\scriptsize\texttt{006832741919}} & M4R & P & \hyperlink{packet-P81}{P81} & Radix-family reassociation with the original coefficient-four resolved source.\\[3pt]
\hypertarget{graph-G02}{\textbf{G02}}\newline{\scriptsize\texttt{01b97022ce29}} & M5 & P & \hyperlink{packet-P81}{P81} & Formula-constructor coefficient five, with explicit logical recoding and physical simulation.\\[3pt]
\hypertarget{graph-G03}{\textbf{G03}}\newline{\scriptsize\texttt{0194cae2b303}} & M6 & P & \hyperlink{packet-P81}{P81} & Formula-constructor coefficient six, with translated numerical calculus and fresh physical source binding.\\[3pt]
\hypertarget{graph-G04}{\textbf{G04}}\newline{\scriptsize\texttt{0049577ad974}} & PA194 (68/8) & P & \hyperlink{packet-P81}{P81} & Same separate-diagonal numerical calculus and original sum guard, alternate 68-reader/eight-fold association.\\[3pt]
\hypertarget{graph-G05}{\textbf{G05}}\newline{\scriptsize\texttt{00564d2d8215}} & PA194-AD & P & \hyperlink{packet-P81}{P81} & Primary 67/7 decoder with A/D physical register placement; source unchanged.\\[3pt]
\hypertarget{graph-G06}{\textbf{G06}}\newline{\scriptsize\texttt{0026de8b5c10}} & PA1944L & P & \hyperlink{packet-P78}{P78} & Four-increment source-tree association ((i,(i,i)),i) with one extra helper and fold;135+67-8 controls; primary source unchanged.\\[3pt]
\hypertarget{graph-G07}{\textbf{G07}}\newline{\scriptsize\texttt{07306c9d5564}} & PA194A & P & \hyperlink{packet-P81}{P81} & Entry jump becomes a decrement of a source-proved zero register; distinct complete initialization.\\[3pt]
\hypertarget{graph-G08}{\textbf{G08}}\newline{\scriptsize\texttt{012d2c9473b0}} & PA194AX11 & P & \hyperlink{packet-P81}{P81} & AX pushes its final guard and falls through into the MP parameter reader; an extra AX operand produces an eleven-entry stack prefix.\\[3pt]
\hypertarget{graph-G09}{\textbf{G09}}\newline{\scriptsize\texttt{002114f775a9}} & PA194AX11H & P & \hyperlink{packet-P79}{P79} & AX11 four-operand/eleven-entry interface composed with cached-head-only H guard and no-push MP.\\[3pt]
\hypertarget{graph-G10}{\textbf{G10}}\newline{\scriptsize\texttt{01792b3f6eb9}} & PA194AX11Q6 & P & \hyperlink{packet-P76}{P76} & AX11 interface, coefficient-six coding, Q head, no-push MP, original sum-based Shift;135+67-8 controls.\\[3pt]
\hypertarget{graph-G11}{\textbf{G11}}\newline{\scriptsize\texttt{0437b6748743}} & PA194AX11QH & P & \hyperlink{packet-P80}{P80} & AX11 interface composed with exchanged Q premises, no-push MP and H guard at coefficient four.\\[3pt]
\hypertarget{graph-G12}{\textbf{G12}}\newline{\scriptsize\texttt{01df63a09c8b}} & PA194AX11QH5 & P & \hyperlink{packet-P82}{P82} & AX11 four-field/eleven-entry interface, coefficient-five coding, Q head, no-push MP, H cache-only Shift; present identically in a filed packet and the latest combined incoming packet.\\[3pt]
\hypertarget{graph-G13}{\textbf{G13}}\newline{\scriptsize\texttt{00c6fd73126c}} & PA194AX11S & P & \hyperlink{packet-P81}{P81} & AX11 interface composed with symmetric original sum-based Shift; coefficient four and original AX formula heads retained.\\[3pt]
\hypertarget{graph-G14}{\textbf{G14}}\newline{\scriptsize\texttt{0081dfa079fd}} & PA194B & P & \hyperlink{packet-P73}{P73} & Four-row history-sensitive physical decoder rewrite: merge reader roles 88/301 while repurposing301 as a dedicated selector; unchanged logical source.\\[3pt]
\hypertarget{graph-G15}{\textbf{G15}}\newline{\scriptsize\texttt{0401717a68a1}} & PA194C & P & \hyperlink{packet-P81}{P81} & Cached-result MP and cached-head-only Shift; duplicated as the principal table by five packets.\\[3pt]
\hypertarget{graph-G16}{\textbf{G16}}\newline{\scriptsize\texttt{0234304c93fc}} & PA194CS & P & \hyperlink{packet-P81}{P81} & Cached-result MP and cached-head-only Shift; symmetric destructive comparison removes a transfer loop.\\[3pt]
\hypertarget{graph-G17}{\textbf{G17}}\newline{\scriptsize\texttt{00877b84b454}} & PA194CSZ = PA194X & P & \hyperlink{packet-P81}{P81} & Cached symmetric controller with known-zero entry decrement; two names identify the same table.\\[3pt]
\hypertarget{graph-G18}{\textbf{G18}}\newline{\scriptsize\texttt{00c1c7c824b6}} & PA194D & P & \hyperlink{packet-P81}{P81} & Separate diagonal head Gamma(A) changes to Gamma(D); revised selectors and 69-reader/nine-fold realization.\\[3pt]
\hypertarget{graph-G19}{\textbf{G19}}\newline{\scriptsize\texttt{013242b84f43}} & PA194DHS & P & \hyperlink{packet-P81}{P81} & Gamma(D) emitter, no-push MP and cached-only symmetric Shift; 135 + 69 - 10 controls.\\[3pt]
\hypertarget{graph-G20}{\textbf{G20}}\newline{\scriptsize\texttt{006e2d8dd8ac}} & PA194F & P & \hyperlink{packet-P81}{P81} & Cached symmetric AX final jump becomes a decrement of X, proved zero by a four-site cut.\\[3pt]
\hypertarget{graph-G21}{\textbf{G21}}\newline{\scriptsize\texttt{047ef5edd6af}} & PA194G & P & \hyperlink{packet-P73}{P73} & AX saved generalization uses its explicit first operand c, giving U(c,h) instead of U(d,h); Q head, no-push MP and sum Shift retained.\\[3pt]
\hypertarget{graph-G22}{\textbf{G22}}\newline{\scriptsize\texttt{0230dd8b1a48}} & PA194H & P & \hyperlink{packet-P81}{P81} & No-push visible-top MP and cached-head-only Shift before symmetric-guard simplification.\\[3pt]
\hypertarget{graph-G23}{\textbf{G23}}\newline{\scriptsize\texttt{0360f9199c19}} & PA194H5 & P & \hyperlink{packet-P81}{P81} & Coefficient-five specialization of the existing H6 command protocol; a new composition, not a new proof-search architecture.\\[3pt]
\hypertarget{graph-G24}{\textbf{G24}}\newline{\scriptsize\texttt{004da46b5d12}} & PA194H6 & P & \hyperlink{packet-P81}{P81} & Coefficient-six formula arithmetic, no-push MP, and cached-only Shift.\\[3pt]
\hypertarget{graph-G25}{\textbf{G25}}\newline{\scriptsize\texttt{007fdfed2daa}} & PA194HA & P & \hyperlink{packet-P81}{P81} & Cached-result MP, original sum Shift, and A/D physical-register placement.\\[3pt]
\hypertarget{graph-G26}{\textbf{G26}}\newline{\scriptsize\texttt{02035799c1ae}} & PA194HAZ & P & \hyperlink{packet-P81}{P81} & HA plus final AX decrement of a register proved zero by a four-site invariant.\\[3pt]
\hypertarget{graph-G27}{\textbf{G27}}\newline{\scriptsize\texttt{0288aecab0a4}} & PA194HS = filed PA194NS & P & \hyperlink{packet-P81}{P81} & No-push visible-top MP combined with cached-head-only symmetric Shift; revised reserved-zero proof compiler.\\[3pt]
\hypertarget{graph-G28}{\textbf{G28}}\newline{\scriptsize\texttt{01926921879c}} & PA194HX & P & \hyperlink{packet-P82}{P82} & Physical H/X register-block exchange; primary resolved source unchanged.\\[3pt]
\hypertarget{graph-G29}{\textbf{G29}}\newline{\scriptsize\texttt{0145cabce701}} & PA194J & P & \hyperlink{packet-P81}{P81} & Ten jumps become known-zero decrements under a finite source invariant; three native rows change.\\[3pt]
\hypertarget{graph-G30}{\textbf{G30}}\newline{\scriptsize\texttt{0130c4d600eb}} & PA194N & P & \hyperlink{packet-P81}{P81} & Successful MP omits dummy-zero push, returns visible consequent (b,w,b,0), and retains sum guard.\\[3pt]
\hypertarget{graph-G31}{\textbf{G31}}\newline{\scriptsize\texttt{003f8e67f486}} & PA194N5 & P & \hyperlink{packet-P81}{P81} & Coefficient-five formula arithmetic plus no-push MP.\\[3pt]
\hypertarget{graph-G32}{\textbf{G32}}\newline{\scriptsize\texttt{007ae5c9ba0a}} & PA194N6 & P & \hyperlink{packet-P81}{P81} & Coefficient-six formula arithmetic plus no-push MP; two packets deliver the same table.\\[3pt]
\hypertarget{graph-G33}{\textbf{G33}}\newline{\scriptsize\texttt{00d2a522634d}} & PA194NS (sum guard) & P & \hyperlink{packet-P81}{P81} & Original sum-testing Shift with symmetric comparison and no-push visible-result MP; distinct from cache-only NS=HS.\\[3pt]
\hypertarget{graph-G34}{\textbf{G34}}\newline{\scriptsize\texttt{00e5a98cd378}} & PA194P5 & P & \hyperlink{packet-P77}{P77} & Physical movement of five registers and seven router targets; logical source unchanged.\\[3pt]
\hypertarget{graph-G35}{\textbf{G35}}\newline{\scriptsize\texttt{00313fc83994}} & PA194Q & P & \hyperlink{packet-P81}{P81} & Exchanges the curried premises in universal distribution, with cached controller; 79 primitive source instructions change.\\[3pt]
\hypertarget{graph-G36}{\textbf{G36}}\newline{\scriptsize\texttt{001f02fcf321}} & PA194Q5 & P & \hyperlink{packet-P81}{P81} & Q emitter plus coefficient-five arithmetic and original dummy-push MP / sum Shift.\\[3pt]
\hypertarget{graph-G37}{\textbf{G37}}\newline{\scriptsize\texttt{0017a41d83cc}} & PA194Q6 & P & \hyperlink{packet-P81}{P81} & Q emitter plus coefficient-six arithmetic and original dummy-push MP / sum Shift.\\[3pt]
\hypertarget{graph-G38}{\textbf{G38}}\newline{\scriptsize\texttt{00600735d0a3}} & PA194QH & P & \hyperlink{packet-P81}{P81} & Q premise exchange with no-push visible-result MP and cache-only Shift.\\[3pt]
\hypertarget{graph-G39}{\textbf{G39}}\newline{\scriptsize\texttt{00b65eb68108}} & PA194QH5 & P & \hyperlink{packet-P81}{P81} & Coefficient five, exchanged universal-distribution premises, visible-result no-push MP, and cache-only Shift.\\[3pt]
\hypertarget{graph-G40}{\textbf{G40}}\newline{\scriptsize\texttt{01ab86d5af64}} & PA194QH6 & P & \hyperlink{packet-P81}{P81} & Q emitter, coefficient-six arithmetic, no-push MP, and cached-only Shift.\\[3pt]
\hypertarget{graph-G41}{\textbf{G41}}\newline{\scriptsize\texttt{0128a632bb8d}} & PA194QHS & P & \hyperlink{packet-P81}{P81} & Exchanged universal-distribution premises, visible no-push MP, and symmetric cache-only Shift, removing a restoration loop relative to QH.\\[3pt]
\hypertarget{graph-G42}{\textbf{G42}}\newline{\scriptsize\texttt{007f19e4bf45}} & PA194QN = PA194NQ & P & \hyperlink{packet-P81}{P81} & Premise-exchanged Q emitter plus no-push N controller and original sum Shift; two deliveries are byte-identical.\\[3pt]
\hypertarget{graph-G43}{\textbf{G43}}\newline{\scriptsize\texttt{0002957d3076}} & PA194QN5 & P & \hyperlink{packet-P81}{P81} & Coefficient five with exchanged universal-distribution premises and no-push MP; the original sum-based Shift remains.\\[3pt]
\hypertarget{graph-G44}{\textbf{G44}}\newline{\scriptsize\texttt{003f0c0e46e0}} & PA194QN6 & P & \hyperlink{packet-P81}{P81} & Coefficient six with exchanged universal-distribution premises and no-push MP; the original sum-based Shift remains.\\[3pt]
\hypertarget{graph-G45}{\textbf{G45}}\newline{\scriptsize\texttt{011894b861f9}} & PA194QNS & P & \hyperlink{packet-P81}{P81} & Q emitter, no-push MP and symmetric original sum Shift; two deliveries are byte-identical.\\[3pt]
\hypertarget{graph-G46}{\textbf{G46}}\newline{\scriptsize\texttt{001967a1e641}} & PA194QS & P & \hyperlink{packet-P81}{P81} & Q emitter, original dummy-push MP and symmetric original sum Shift.\\[3pt]
\hypertarget{graph-G47}{\textbf{G47}}\newline{\scriptsize\texttt{00a1583fcbbb}} & PA194R & P & \hyperlink{packet-P81}{P81} & F zero-closing source with reassociated 67-reader grammar; 1111 source addresses change.\\[3pt]
\hypertarget{graph-G48}{\textbf{G48}}\newline{\scriptsize\texttt{0096f438b79b}} & PA194SUMS & P & \hyperlink{packet-P81}{P81} & Original sum-testing Shift with symmetric comparison, preserving primary PA194 command semantics.\\[3pt]
\hypertarget{graph-G49}{\textbf{G49}}\newline{\scriptsize\texttt{000015e64f03}} & PA194TDY & P & \hyperlink{packet-P75}{P75} & Physical T/D/Y register-cycle allocation; logical source unchanged.\\[3pt]
\hypertarget{graph-G50}{\textbf{G50}}\newline{\scriptsize\texttt{004bf9b352ca}} & PA194ZR & P & \hyperlink{packet-P81}{P81} & Eight comparator returns and AX exit use zero-decrement return; independent 129-site invariant.\\[3pt]
\hypertarget{graph-G51}{\textbf{G51}}\newline{\scriptsize\texttt{002c47745da7}} & Primary PA194 (67/7) & P & \hyperlink{packet-P80}{P80} & Separate diagonal emitter, original sum-guard controller; 134 fixed + 67 readers - 7 folds.\\[3pt]
\hypertarget{graph-G52}{\textbf{G52}}\newline{\scriptsize\texttt{000adfa7c076}} & C-I-Z & D & \hyperlink{packet-P81}{P81} & Auxiliary cached-source/zero-return search input; no principal release claim.\\[3pt]
\hypertarget{graph-G53}{\textbf{G53}}\newline{\scriptsize\texttt{02c0dea52e0b}} & CS both-zero6 & D & \hyperlink{packet-P81}{P81} & Auxiliary simultaneous zero-return search input; no principal consistency proof claim.\\[3pt]
\hypertarget{graph-G54}{\textbf{G54}}\newline{\scriptsize\texttt{006b2495744f}} & CSZA final-register 3 & S & \hyperlink{packet-P81}{P81} & Known-zero final-return sibling with construction/startup evidence.\\[3pt]
\hypertarget{graph-G55}{\textbf{G55}}\newline{\scriptsize\texttt{00c989810d4b}} & CSZA final-register 5 & S & \hyperlink{packet-P81}{P81} & Known-zero final-return sibling; not included among independently replayed main D/J machines.\\[3pt]
\hypertarget{graph-G56}{\textbf{G56}}\newline{\scriptsize\texttt{01195131ce89}} & CSZA final-register 6 & S & \hyperlink{packet-P81}{P81} & Known-zero final-return sibling with construction/startup evidence.\\[3pt]
\hypertarget{graph-G57}{\textbf{G57}}\newline{\scriptsize\texttt{01cc2aa9d7aa}} & Gamma(D) + CSZA & S & \hyperlink{packet-P81}{P81} & Gamma(D) emitter with a known-zero cached symmetric controller; only supplementary evidence claimed.\\[3pt]
\hypertarget{graph-G58}{\textbf{G58}}\newline{\scriptsize\texttt{0100da703594}} & MPH auxiliary & D & \hyperlink{packet-P81}{P81} & MP returns its result in the cache while retaining the original sum guard; incoming HAZ credits this prior table.\\[3pt]
\hypertarget{graph-G59}{\textbf{G59}}\newline{\scriptsize\texttt{038656d61b03}} & N6 append-zero tail & R & \hyperlink{packet-P81}{P81} & Unrepaired append-one to append-zero transplantation; no calibrated blank startup.\\[3pt]
\hypertarget{graph-G60}{\textbf{G60}}\newline{\scriptsize\texttt{0162e24e5c9f}} & PA194-altC & D & \hyperlink{packet-P81}{P81} & Alternate cached decoder in grammar search; not a separate principal verified theorem.\\[3pt]
\hypertarget{graph-G61}{\textbf{G61}}\newline{\scriptsize\texttt{04e51ef4bb79}} & PA194CZ & S & \hyperlink{packet-P81}{P81} & Cached controller plus known-zero source jump; a secondary same-size composition.\\[3pt]
\hypertarget{graph-G62}{\textbf{G62}}\newline{\scriptsize\texttt{00ea59f8ca73}} & PA194F\_AD & S & \hyperlink{packet-P81}{P81} & A/D register placement composed with known-zero final emitter return; supplied physical/source checks.\\[3pt]
\hypertarget{graph-G63}{\textbf{G63}}\newline{\scriptsize\texttt{0137737f5a6a}} & PA194N-AD (auxiliary) & A & \hyperlink{packet-P81}{P81} & No-push N plus A/D physical-register placement; originally JSON-only comparison candidate, now retained as a literal table.\\[3pt]
\hypertarget{graph-G64}{\textbf{G64}}\newline{\scriptsize\texttt{008b9d3f990e}} & PA194R\_AD & S & \hyperlink{packet-P81}{P81} & A/D placement composed with reassociated decoder and known-zero final emitter return; supplied physical/source checks.\\[3pt]
\hypertarget{graph-G65}{\textbf{G65}}\newline{\scriptsize\texttt{019c14cb5a0a}} & Q with original controller (QS parent) & A & \hyperlink{packet-P81}{P81} & Q emitter with original dummy-push MP and original sum Shift; retained parent of the QS construction, not a separate principal theorem.\\[3pt]
\hypertarget{graph-G66}{\textbf{G66}}\newline{\scriptsize\texttt{0033bf5e00cc}} & Rejected append-zero tail & R & \hyperlink{packet-P81}{P81} & Literal append-one to append-zero transplant fails its blank-start behavior: all 194 starts classified, four halt and 190 enter proven blank rays.\\[3pt]
\hypertarget{graph-G67}{\textbf{G67}}\newline{\scriptsize\texttt{0076a77b8128}} & Rejected comparator return 103 & R & \hyperlink{packet-P81}{P81} & Source-equivalent patch passes source contracts but halts after 1361 blank-start steps.\\[3pt]
\hypertarget{graph-G68}{\textbf{G68}}\newline{\scriptsize\texttt{0022fe9abb7a}} & Rejected dec6 & R & \hyperlink{packet-P81}{P81} & Rejected structural return/decrement candidate; do not count as a consistency construction.\\[3pt]
\hypertarget{graph-G69}{\textbf{G69}}\newline{\scriptsize\texttt{0156af562b93}} & Rejected even-runtime194 & R & \hyperlink{packet-P76}{P76} & Even scheduler combined with Runtime193 quotient; startup enters a certified outward ray before reaching new scheduler control.\\[3pt]
\hypertarget{graph-G70}{\textbf{G70}}\newline{\scriptsize\texttt{00002feec556}} & Rejected node1-X & R & \hyperlink{packet-P81}{P81} & Known-zero-looking patch is rejected; source correctness alone does not ensure blank initialization.\\[3pt]
\hypertarget{graph-G71}{\textbf{G71}}\newline{\scriptsize\texttt{0054b7f16539}} & Rejected short scanner & R & \hyperlink{packet-P81}{P81} & Seven-zero scanner quotient violates a legal source contract; not an accepted consistency machine.\\[3pt]
\hypertarget{graph-G72}{\textbf{G72}}\newline{\scriptsize\texttt{01772d2f9520}} & Source rotation 100\_164\_left & A & \hyperlink{packet-P81}{P81} & Physical source-tree rotation included as search input; no separate complete consistency proof attributed.\\[3pt]
\hypertarget{graph-G73}{\textbf{G73}}\newline{\scriptsize\texttt{037f9362af4d}} & Source rotation 100\_174\_right & A & \hyperlink{packet-P81}{P81} & Physical source-tree rotation included as search input; no separate complete consistency proof attributed.\\[3pt]
\hypertarget{graph-G74}{\textbf{G74}}\newline{\scriptsize\texttt{004774a044e6}} & Source rotation 166\_179\_left & A & \hyperlink{packet-P81}{P81} & Physical source-tree rotation included as search input; no separate complete consistency proof attributed.\\[3pt]
\hypertarget{graph-G75}{\textbf{G75}}\newline{\scriptsize\texttt{007d8226ee9e}} & Source rotation 19\_29\_left & A & \hyperlink{packet-P81}{P81} & Physical source-tree rotation included as search input; no separate complete consistency proof attributed.\\[3pt]
\hypertarget{graph-G76}{\textbf{G76}}\newline{\scriptsize\texttt{000e4f053321}} & Source rotation 2000\_2001\_right & A & \hyperlink{packet-P81}{P81} & Physical source-tree rotation included as search input; no separate complete consistency proof attributed.\\[3pt]
\hypertarget{graph-G77}{\textbf{G77}}\newline{\scriptsize\texttt{00c8200ecb4f}} & Source rotation 2007\_2008\_left & A & \hyperlink{packet-P81}{P81} & Physical source-tree rotation included as search input; no separate complete consistency proof attributed.\\[3pt]
\hypertarget{graph-G78}{\textbf{G78}}\newline{\scriptsize\texttt{0201d85b3a8b}} & Source rotation 2008\_2009\_right & A & \hyperlink{packet-P81}{P81} & Physical source-tree rotation included as search input; no separate complete consistency proof attributed.\\[3pt]
\hypertarget{graph-G79}{\textbf{G79}}\newline{\scriptsize\texttt{01662a3af6a6}} & Source rotation 2009\_2010\_right & A & \hyperlink{packet-P81}{P81} & Physical source-tree rotation included as search input; no separate complete consistency proof attributed.\\[3pt]
\hypertarget{graph-G80}{\textbf{G80}}\newline{\scriptsize\texttt{0002b0f09c18}} & Source rotation 2009\_2011\_left & A & \hyperlink{packet-P81}{P81} & Physical source-tree rotation included as search input; no separate complete consistency proof attributed.\\[3pt]
\hypertarget{graph-G81}{\textbf{G81}}\newline{\scriptsize\texttt{0083f916c7a0}} & Source rotation 2011\_2012\_left & A & \hyperlink{packet-P81}{P81} & Physical source-tree rotation included as search input; no separate complete consistency proof attributed.\\[3pt]
\hypertarget{graph-G82}{\textbf{G82}}\newline{\scriptsize\texttt{004d8495d142}} & Source rotation 2040\_2041\_right & A & \hyperlink{packet-P81}{P81} & Physical source-tree rotation included as search input; no separate complete consistency proof attributed.\\[3pt]
\hypertarget{graph-G83}{\textbf{G83}}\newline{\scriptsize\texttt{00ed0dd04b89}} & Source rotation 2041\_2042\_right & A & \hyperlink{packet-P81}{P81} & Physical source-tree rotation included as search input; no separate complete consistency proof attributed.\\[3pt]
\hypertarget{graph-G84}{\textbf{G84}}\newline{\scriptsize\texttt{01037b5931e9}} & Source rotation 2041\_2047\_left & A & \hyperlink{packet-P81}{P81} & Physical source-tree rotation included as search input; no separate complete consistency proof attributed.\\[3pt]
\hypertarget{graph-G85}{\textbf{G85}}\newline{\scriptsize\texttt{015d504d9c5e}} & Source rotation 2047\_2048\_right & A & \hyperlink{packet-P81}{P81} & Physical source-tree rotation included as search input; no separate complete consistency proof attributed.\\[3pt]
\hypertarget{graph-G86}{\textbf{G86}}\newline{\scriptsize\texttt{0023aff57dc9}} & Source rotation 2047\_2053\_left & A & \hyperlink{packet-P81}{P81} & Physical source-tree rotation included as search input; no separate complete consistency proof attributed.\\[3pt]
\hypertarget{graph-G87}{\textbf{G87}}\newline{\scriptsize\texttt{0109a111bb4d}} & Source rotation 2048\_2049\_left & A & \hyperlink{packet-P81}{P81} & Physical source-tree rotation included as search input; no separate complete consistency proof attributed.\\[3pt]
\hypertarget{graph-G88}{\textbf{G88}}\newline{\scriptsize\texttt{02552a526434}} & Source rotation 2053\_2054\_right & A & \hyperlink{packet-P81}{P81} & Physical source-tree rotation included as search input; no separate complete consistency proof attributed.\\[3pt]
\hypertarget{graph-G89}{\textbf{G89}}\newline{\scriptsize\texttt{028891e49882}} & Source rotation 2061\_2062\_left & A & \hyperlink{packet-P81}{P81} & Physical source-tree rotation included as search input; no separate complete consistency proof attributed.\\[3pt]
\hypertarget{graph-G90}{\textbf{G90}}\newline{\scriptsize\texttt{035c86205af5}} & Source rotation 214\_312\_left & A & \hyperlink{packet-P81}{P81} & Physical source-tree rotation included as search input; no separate complete consistency proof attributed.\\[3pt]
\hypertarget{graph-G91}{\textbf{G91}}\newline{\scriptsize\texttt{004dbf2558b5}} & Source rotation 217\_197\_right & A & \hyperlink{packet-P81}{P81} & Physical source-tree rotation included as search input; no separate complete consistency proof attributed.\\[3pt]
\hypertarget{graph-G92}{\textbf{G92}}\newline{\scriptsize\texttt{001e37fdea79}} & Source rotation 221\_217\_left & A & \hyperlink{packet-P81}{P81} & Physical source-tree rotation included as search input; no separate complete consistency proof attributed.\\[3pt]
\hypertarget{graph-G93}{\textbf{G93}}\newline{\scriptsize\texttt{0230adf3c030}} & Source rotation 29\_44\_left & A & \hyperlink{packet-P81}{P81} & Physical source-tree rotation included as search input; no separate complete consistency proof attributed.\\[3pt]
\hypertarget{graph-G94}{\textbf{G94}}\newline{\scriptsize\texttt{0166cb1df1f1}} & Source rotation 312\_38\_left & A & \hyperlink{packet-P81}{P81} & Physical source-tree rotation included as search input; no separate complete consistency proof attributed.\\[3pt]
\end{longtable}
\endgroup

\clearpage
\section{The 82-packet delivery ledger}

\label{app:packet-ledger}

This ledger preserves distinct deliveries even when their principal tables coincide. The final column summarizes the packet's own reported contribution; historical statements about PA proofs or calibration are attributed claims, not newly discharged iPRA endpoints. In particular, a packet's former state record reflects its date. The five incoming packets at the initial snapshot are P78--P82. All 84 top-level README files in these 82 packets were located and hashed.

For P01--P77, the exact repository directory is \nolinkurl{Computability/BusyBeaver/Research/PA/docs/} followed by the packet name. For P78--P82, the exact repository file is \nolinkurl{docs/incoming/} followed by the packet name and \nolinkurl{.zip}. The snapshot is \texttt{32d0d8b7d93a63ad4106c18cdb29b37927acdad2}. Full source and README identities are retained in \nolinkurl{comparison/packet_ledger.csv} and \nolinkurl{comparison/packet_ledger.json}.
\begingroup
\footnotesize
\setlength{\tabcolsep}{2pt}
\setlength{\LTleft}{0pt}
\setlength{\LTright}{0pt}
\renewcommand{\arraystretch}{1.12}
\begin{longtable}{@{}>{\raggedright\arraybackslash}p{.055\linewidth}>{\raggedright\arraybackslash}p{.355\linewidth}>{\raggedright\arraybackslash}p{.55\linewidth}@{}}
\caption{All filed and incoming packets at the initial snapshot}\label{tab:packet-ledger}\\
\toprule
ID & Packet name & Reported contribution\\
\midrule
\endfirsthead
\multicolumn{3}{@{}l}{\textit{Table \thetable\ continued}}\\
\toprule
ID & Packet name & Reported contribution\\
\midrule
\endhead
\midrule
\multicolumn{3}{r@{}}{\textit{Continued on the next page}}\\
\endfoot
\bottomrule
\endlastfoot
\hypertarget{packet-P01}{\textbf{P01}} & \nolinkurl{pa_unary_power_cover} & Unary power covering at PA strength (report of 27 October 2026). The U208 packets carry it as a baseline.\\[3pt]
\hypertarget{packet-P02}{\textbf{P02}} & \nolinkurl{U208_ConPA_Conditional_Context_2026-10-29} & U208 by a conditional tape-context reduction. States a 208-state candidate \nolinkurl{machine/u208.tm}.\\[3pt]
\hypertarget{packet-P03}{\textbf{P03}} & \nolinkurl{U208_ConPA_Entry_Context_2026-10-29} & U208 by eliminating a state using entry context. States a 208-state machine \nolinkurl{machine/u208.tm}.\\[3pt]
\hypertarget{packet-P04}{\textbf{P04}} & \nolinkurl{U208_ConPA_History_2026-10-29} & U208 by a finite-history decoder fold. States a 208-state candidate \nolinkurl{machine/u208.tm}.\\[3pt]
\hypertarget{packet-P05}{\textbf{P05}} & \nolinkurl{U208_ConPA_2026-10-29} & U208 by a three-cell fold. States a 208-state candidate \nolinkurl{machine/u208.tm}.\\[3pt]
\hypertarget{packet-P06}{\textbf{P06}} & \nolinkurl{u208_research_2026-10-29} & U208 research release. States two 208-state machines, \nolinkurl{machine/u208.tm} and \nolinkurl{machine/variant_plain/u208.tm}.\\[3pt]
\hypertarget{packet-P07}{\textbf{P07}} & \nolinkurl{Below_U208_Research_2026-10-29} & Research toward fewer than 208 states. States that no machine with 207 states or fewer was obtained.\\[3pt]
\hypertarget{packet-P08}{\textbf{P08}} & \nolinkurl{u208_comparison_research_2026-10-29} & Comparison of the 208-state proposals and a search below 208. States that the target was not achieved.\\[3pt]
\hypertarget{packet-P09}{\textbf{P09}} & \nolinkurl{ConPA_below208_research_2026-10-29} & Comparative audit and reduction search below 208. States that it does not establish a 207-state machine.\\[3pt]
\hypertarget{packet-P10}{\textbf{P10}} & \nolinkurl{pa_consistency_research_2026-10-29} & Audit and new source reductions. States that no validated machine below 208 states was obtained.\\[3pt]
\hypertarget{packet-P11}{\textbf{P11}} & \nolinkurl{pa208_research_2026-10-29} & Comparison and research below 208 states. States that it does not establish a valid machine below 208.\\[3pt]
\hypertarget{packet-P12}{\textbf{P12}} & \nolinkurl{pa_tm_research_2026-11-13} & Comparison and reduction research. States a new 196-state variant with shorter source paths and nothing below 196.\\[3pt]
\hypertarget{packet-P13}{\textbf{P13}} & \nolinkurl{pa_consistency_research_2026-11-13} & Comparison and diagonal-vacuity research. States a shorter axiom program with a 199-state implementation and nothing below 196.\\[3pt]
\hypertarget{packet-P14}{\textbf{P14}} & \nolinkurl{conpa_tm_comparison_and_reduction_search} & Comparison and reduction search. States two experimental 196-state command variants and nothing below 196.\\[3pt]
\hypertarget{packet-P15}{\textbf{P15}} & \nolinkurl{pa196_research_2026-11-13} & PA196 comparison, source factorization and proof obligations. States a source tree of 157 internal nodes and nothing below 196.\\[3pt]
\hypertarget{packet-P16}{\textbf{P16}} & \nolinkurl{pa_consistency_research} & Research and verification packet. States that no machine below 196 states has been validated.\\[3pt]
\hypertarget{packet-P17}{\textbf{P17}} & \nolinkurl{pa_tm_comparison_and_reductions} & Comparison and new rewrites, a research snapshot.\\[3pt]
\hypertarget{packet-P18}{\textbf{P18}} & \nolinkurl{pa_tm_extended_research_2026-11-13} & Comparison and refined reduction search. States that the smallest accepted count is 196.\\[3pt]
\hypertarget{packet-P19}{\textbf{P19}} & \nolinkurl{conpa_tm_comparison_and_search} & Comparison and a search below 196. States that none was obtained; files named \nolinkurl{195} are rejected or diagnostic candidates.\\[3pt]
\hypertarget{packet-P20}{\textbf{P20}} & \nolinkurl{pa_consistency_reduction} & Comparison and reduction experiments. States two new 196-state variants and a shorter source program with a 203-state implementation.\\[3pt]
\hypertarget{packet-P21}{\textbf{P21}} & \nolinkurl{pa_consistency_comparison_2026_11_13} & Comparison and new reduction experiments. States that nothing below 196 was obtained.\\[3pt]
\hypertarget{packet-P22}{\textbf{P22}} & \nolinkurl{PA_consistency_TM_research_2026-11-13} & Comparison and new source reductions. States that no machine below 196 is claimed.\\[3pt]
\hypertarget{packet-P23}{\textbf{P23}} & \nolinkurl{pa_consistency_research_2026-11-13_2} & Comparison and search below 196. States two new layouts with 196 and 198 states and a command-level reduction.\\[3pt]
\hypertarget{packet-P24}{\textbf{P24}} & \nolinkurl{conpa_tm_study_2026-11-13} & Comparison and reduction study. States that nothing below 196 was obtained.\\[3pt]
\hypertarget{packet-P25}{\textbf{P25}} & \nolinkurl{pa_below196_research} & Comparison and reduction experiments. States that nothing below 196 was obtained.\\[3pt]
\hypertarget{packet-P26}{\textbf{P26}} & \nolinkurl{pa_consistency_comparison_and_new196_2026-11-13} & Comparison and reduction research. States that it contains no accepted machine below 196.\\[3pt]
\hypertarget{packet-P27}{\textbf{P27}} & \nolinkurl{pa_tm_research_2026-11-13_2} & Comparison and reduction investigation. States same-size work, failed candidates and nothing below 196.\\[3pt]
\hypertarget{packet-P28}{\textbf{P28}} & \nolinkurl{pa_consistency_cross_archive} & Cross-archive constructions and reduction search. States that it does not certify a 195-state machine.\\[3pt]
\hypertarget{packet-P29}{\textbf{P29}} & \nolinkurl{pa195_conpa_2026-11-13} & PA195. States a 195-state machine \nolinkurl{pa195/pa195.tm} and an article with a PA proof of the equivalence.\\[3pt]
\hypertarget{packet-P30}{\textbf{P30}} & \nolinkurl{pa_consistency_reduction_audit_2026-11-14} & Comparison and reduction audit. States nothing below 196 and no new theorem.\\[3pt]
\hypertarget{packet-P31}{\textbf{P31}} & \nolinkurl{pa_consistency_research_2026-11-14} & Comparison and NEG18. States an 18-instruction encoded negation with a compiled machine of 197 states.\\[3pt]
\hypertarget{packet-P32}{\textbf{P32}} & \nolinkurl{PA195_ConPA_research_2026-11-14} & A 195-state machine. States the primary \nolinkurl{machine_195S.tm} and a second table \nolinkurl{machine_195C.tm}.\\[3pt]
\hypertarget{packet-P33}{\textbf{P33}} & \nolinkurl{pa195_followup_research} & Comparison and searches below 195. States a graph-distinct 195-state machine M195S8 and nothing below 195.\\[3pt]
\hypertarget{packet-P34}{\textbf{P34}} & \nolinkurl{pa194_conpa_2026-11-13} & A 194-state machine. States the table \nolinkurl{construction/tables/pa194.tm} and a proof of the equivalence in PA.\\[3pt]
\hypertarget{packet-P35}{\textbf{P35}} & \nolinkurl{pa194_conpa_2026-11-13_2} & PA194. States a 194-state machine \nolinkurl{pa194/pa194.tm} and an article with a PA proof of the equivalence.\\[3pt]
\hypertarget{packet-P36}{\textbf{P36}} & \nolinkurl{pa195_diagonal_combined_2026-11-13} & PA195DC. States a different 195-state machine.\\[3pt]
\hypertarget{packet-P37}{\textbf{P37}} & \nolinkurl{pa194C_conpa_2026-11-13_2} & PA194C. States a materially different 194-state machine, an article with a PA proof of the equivalence, and nothing below 194.\\[3pt]
\hypertarget{packet-P38}{\textbf{P38}} & \nolinkurl{pa194_new_designs_2026-11-13} & Two new 194-state machines, PA194C and PA194CS. States that nothing below 194 was obtained.\\[3pt]
\hypertarget{packet-P39}{\textbf{P39}} & \nolinkurl{pa194c_conpa} & PA194C, a cached-head 194-state machine. States that no machine below 194 was established.\\[3pt]
\hypertarget{packet-P40}{\textbf{P40}} & \nolinkurl{pa194c_conpa_2026-11-13_1} & PA194C. States a different 194-state machine and an article with a PA proof of the equivalence.\\[3pt]
\hypertarget{packet-P41}{\textbf{P41}} & \nolinkurl{pa194c_conpa_2026-11-13} & PA194C, a cached-head 194-state machine. States an article with a PA proof of the equivalence and nothing below 194.\\[3pt]
\hypertarget{packet-P42}{\textbf{P42}} & \nolinkurl{pa194_radix_family_2026-11-13} & New 194-state machines calibrated to Con(PA), a radix family.\\[3pt]
\hypertarget{packet-P43}{\textbf{P43}} & \nolinkurl{pa194j_pa194d_conpa_2026-11-13} & PA194J and PA194D. States two new 194-state machines.\\[3pt]
\hypertarget{packet-P44}{\textbf{P44}} & \nolinkurl{pa194q_conpa_2026-11-13} & PA194Q and PA194CSZ. States two new 194-state machines and that no machine below 194 was established.\\[3pt]
\hypertarget{packet-P45}{\textbf{P45}} & \nolinkurl{PA194R_research_2026-11-14} & PA194R. States a different 194-state machine with a new decoder layout and that no machine below 194 was verified.\\[3pt]
\hypertarget{packet-P46}{\textbf{P46}} & \nolinkurl{pa194N_conpa_2026-11-13} & PA194N. States a new 194-state machine, an article with a PA proof of the equivalence, and nothing below 194.\\[3pt]
\hypertarget{packet-P47}{\textbf{P47}} & \nolinkurl{pa194zr_conpa_2026-11-13} & PA194ZR. States a new 194-state machine \nolinkurl{machine/pa194zr.tm}, a PA proof of the equivalence, and nothing below 194.\\[3pt]
\hypertarget{packet-P48}{\textbf{P48}} & \nolinkurl{pa194hs_conpa_2026-11-13} & PA194HS and PA194H. States two 194-state machines and that no machine below 194 was established.\\[3pt]
\hypertarget{packet-P49}{\textbf{P49}} & \nolinkurl{pa194n5_n6_conpa_2026-11-13} & PA194N5 and PA194N6. States two distinct 194-state machines and claims no machine below 194.\\[3pt]
\hypertarget{packet-P50}{\textbf{P50}} & \nolinkurl{pa194ns_conpa_2026-11-13} & PA194NS. States a different 194-state machine and that no machine below 194 was established.\\[3pt]
\hypertarget{packet-P51}{\textbf{P51}} & \nolinkurl{pa194n6_research_2026-11-13} & PA194N6. States a different 194-state machine and that no machine below 194 was verified.\\[3pt]
\hypertarget{packet-P52}{\textbf{P52}} & \nolinkurl{PA194NS_research} & PA194NS and PA194SUMS. States two new 194-state constructions.\\[3pt]
\hypertarget{packet-P53}{\textbf{P53}} & \nolinkurl{pa194haz_research_2026-11-14} & PA194HA and PA194HAZ. States two 194-state alternatives and that no machine below 194 was obtained.\\[3pt]
\hypertarget{packet-P54}{\textbf{P54}} & \nolinkurl{pa194qn_research_2026-11-13} & PA194QN. States a distinct 194-state machine and that no verified machine below 194 was found.\\[3pt]
\hypertarget{packet-P55}{\textbf{P55}} & \nolinkurl{pa194h6_conpa_2026-11-13} & PA194H6. States a new 194-state machine, an article with a PA proof of the equivalence, and nothing below 194.\\[3pt]
\hypertarget{packet-P56}{\textbf{P56}} & \nolinkurl{pa194dhs_conpa_2026-11-13} & PA194DHS. States a different 194-state machine and that no machine below 194 was verified.\\[3pt]
\hypertarget{packet-P57}{\textbf{P57}} & \nolinkurl{pa194nq_research_2026-11-13} & PA194NQ. States a new 194-state machine and claims no machine below 194.\\[3pt]
\hypertarget{packet-P58}{\textbf{P58}} & \nolinkurl{pa194_qns_research_2026-11-13} & PA194QNS. States a new 194-state machine and that no verified construction below 194 was found.\\[3pt]
\hypertarget{packet-P59}{\textbf{P59}} & \nolinkurl{pa194_qradix_research_2026-11-13} & PA194Q5 and PA194Q6. States two different 194-state machines and that nothing below 194 was obtained.\\[3pt]
\hypertarget{packet-P60}{\textbf{P60}} & \nolinkurl{pa194qh6_conpa_2026-11-13} & PA194QH6. States a new 194-state composition and that no machine below 194 was obtained.\\[3pt]
\hypertarget{packet-P61}{\textbf{P61}} & \nolinkurl{pa194qns_research_2026-11-13} & PA194QNS. States a distinct 194-state machine and that no machine below 194 was established.\\[3pt]
\hypertarget{packet-P62}{\textbf{P62}} & \nolinkurl{pa194qh_conpa_2026-11-13} & PA194QH. States a new 194-state machine combining the Q emitter with the H controller, and that the reduction below 194 was not achieved.\\[3pt]
\hypertarget{packet-P63}{\textbf{P63}} & \nolinkurl{pa_consistency_reduction_study_2026-11-13} & Comparison and reduction study. States that no valid machine below 194 was established and an obstruction for one family of transformations.\\[3pt]
\hypertarget{packet-P64}{\textbf{P64}} & \nolinkurl{pa_consistency_research_2026-11-14_2} & Comparison and investigation below 194 states. States that no verified machine below 194 was obtained, 49 distinct 194-state graph classes and an all-offset obstruction.\\[3pt]
\hypertarget{packet-P65}{\textbf{P65}} & \nolinkurl{pa_sub194_research} & Search below 194 states. States that the bound was not achieved and a new 196-state construction MP196.\\[3pt]
\hypertarget{packet-P66}{\textbf{P66}} & \nolinkurl{pa_sub194_research_2026-11-13} & Comparison and research below 194. States that no machine below 194 has been established and an explicit 193-state runtime core whose tested startups are rejected.\\[3pt]
\hypertarget{packet-P67}{\textbf{P67}} & \nolinkurl{pa194_ax11_research_2026-11-14} & PA194AX11 and PA194QH5. States a different 194-state machine with a PA proof of the equivalence, a second 194-state composition, and nothing below 194.\\[3pt]
\hypertarget{packet-P68}{\textbf{P68}} & \nolinkurl{pa194_qh5_research_2026-11-14} & PA194QH5. States a new 194-state construction and that no machine below 194 is established.\\[3pt]
\hypertarget{packet-P69}{\textbf{P69}} & \nolinkurl{pa194qhs_conpa_2026-11-14} & PA194QHS. States a 194-state composition not delivered before and that no machine below 194 was established.\\[3pt]
\hypertarget{packet-P70}{\textbf{P70}} & \nolinkurl{pa194qn6_conpa_2026-11-14} & PA194QN6. States a new 194-state composition and claims no machine below 194.\\[3pt]
\hypertarget{packet-P71}{\textbf{P71}} & \nolinkurl{pa194qn_2026-11-14} & PA194QN5 and PA194QN6. States two 194-state machines and that no valid machine below 194 was established.\\[3pt]
\hypertarget{packet-P72}{\textbf{P72}} & \nolinkurl{pa194h5_research_2026-11-14} & PA194H5. States a new 194-state coefficient-five specialization of the H architecture and nothing below 194.\\[3pt]
\hypertarget{packet-P73}{\textbf{P73}} & \nolinkurl{PA194B_PA194G_research_2026-11-14} & PA194B and PA194G. States two new 194-state machines and that no verified machine below 194 was obtained.\\[3pt]
\hypertarget{packet-P74}{\textbf{P74}} & \nolinkurl{pa194_ax11qh5_2026-11-14} & PA194AX11QH5. States a new 194-state machine combining an eleven-entry AX protocol with coefficient-five coding, and nothing below 194.\\[3pt]
\hypertarget{packet-P75}{\textbf{P75}} & \nolinkurl{pa194_tdy_2026-11-14} & PA194TDY. States a 194-state physical register-cycle design and that no machine below 194 was established.\\[3pt]
\hypertarget{packet-P76}{\textbf{P76}} & \nolinkurl{pa194ax11q6_2026-11-14} & PA194AX11Q6. States a new 194-state composition and claims neither a machine below 194 nor a lower bound.\\[3pt]
\hypertarget{packet-P77}{\textbf{P77}} & \nolinkurl{pa194p5_research_2026-11-14} & PA194P5. States a 194-state machine that moves five registers and seven router targets of the primary PA194, a PA proof that it halts exactly when PA194 does, and nothing below 194.\\[3pt]
\hypertarget{packet-P78}{\textbf{P78}} & \nolinkurl{pa1944l_research_2026-11-14} & PA1944L: new physical four-increment decoder, 135 fixed + 67 AX readers - 8 folds = 194; same primary source. Written relative iPRA halting equivalence; absolute iPRA calibration explicitly unproved.\\[3pt]
\hypertarget{packet-P79}{\textbf{P79}} & \nolinkurl{pa194_ax11h_research_2026-11-14} & PA194AX11H: AX11 interface with H cache-only guard; 194 controls. Corrected relative refutation compiler and literal-tape certificates; absolute iPRA calibration explicitly unproved.\\[3pt]
\hypertarget{packet-P80}{\textbf{P80}} & \nolinkurl{pa194_ax11qh_research_2026-11-14} & PA194AX11QH: AX11, exchanged Q premises, no-push MP, H guard; 194 controls. Written relative iPRA halting equivalence; absolute calibration explicitly unproved.\\[3pt]
\hypertarget{packet-P81}{\textbf{P81}} & \nolinkurl{pa194_ax11s_study} & PA194AX11S: AX11 plus symmetric original sum guard; 194 controls. Written relative source-equivalence proof; absolute iPRA calibration explicitly unproved.\\[3pt]
\hypertarget{packet-P82}{\textbf{P82}} & \nolinkurl{pa194_research_2026-11-14} & AX11QH5 and PA194HX: coefficient-five AX11/Q/H composition and primary physical H/X swap, both 194. Command-level weak-arithmetic proof and tape certificates; absolute iPRA calibration explicitly unproved.\\[3pt]
\end{longtable}
\endgroup

\end{document}
